% BEGIN EMBEDDED BIBLIOGRAPHY
% Bibliography data are part of this source; the compiler writes its temporary input.
\begin{filecontents*}[overwrite]{\jobname.bib}
@unpublished{QED3Manuscript,
  title = {The field-strength {Osterwalder--Schrader} axioms for {QED} in three dimensions},
  date = {2026},
  shorthand = {Cur},
  note = {Comprehensive edition, 263 pp., 20 October 2026, before the modular integration revision}
}
\end{filecontents*}
% END EMBEDDED BIBLIOGRAPHY
\documentclass[11pt,letterpaper,reqno]{amsart}
% BEGIN INLINE manuscript-style
\makeatletter
\NeedsTeXFormat{LaTeX2e}[2020/10/01]
% Use with \documentclass[11pt,letterpaper,reqno]{amsart}.
% This file controls presentation only. It supplies no mathematical content.

\RequirePackage[T1]{fontenc}
\RequirePackage{lmodern}
\RequirePackage{microtype}
\RequirePackage[letterpaper,margin=1in,headheight=15pt]{geometry}
\RequirePackage{amsmath,amsthm,amssymb,mathtools}
\RequirePackage{enumitem}
\RequirePackage{perpage}
\RequirePackage{xcolor}
\definecolor{manuscriptred}{rgb}{0.6,0.1,0.1}
\definecolor{manuscriptblue}{rgb}{0.1,0.3,0.5}
\definecolor{manuscriptmagenta}{HTML}{60445C}
\RequirePackage[breaklinks=true,colorlinks=true,
  linkcolor=manuscriptred,citecolor=manuscriptblue,urlcolor=manuscriptmagenta]{hyperref}
\RequirePackage{cleveref}
\hypersetup{pdfauthor={}}

% Retain the AMS paragraph, title, heading and running-head conventions.
% Equations retain one global Arabic sequence, as in the reference.
\setcounter{tocdepth}{1}
\renewcommand{\thefigure}{\Roman{figure}}
\renewcommand{\thetable}{\Roman{table}}
\MakePerPage{footnote}
\renewcommand{\thefootnote}{%
  \ifnum\value{footnote}>9\arabic{footnote}\else\fnsymbol{footnote}\fi}

% Upright statements. Main results are A, B, C; supporting statements
% share a section-local sequence. These end marks indicate typography,
% never proof completeness, validation, or a deferred-proof status.
\theoremstyle{definition}
\newtheorem{ms@theorem}{Theorem}
\renewcommand{\thems@theorem}{\Alph{ms@theorem}}
\newtheorem*{ms@theoremstar}{Theorem}
\newtheorem{ms@proposition}{Proposition}[section]
\newtheorem{ms@lemma}[ms@proposition]{Lemma}
\newtheorem{ms@definition}[ms@proposition]{Definition}
\newtheorem{ms@corollary}[ms@proposition]{Corollary}
\newtheorem*{ms@corollarystar}{Corollary}

\newcommand{\ms@majorstart}{%
  \pushQED{\qed}\renewcommand{\qedsymbol}{$\blacksquare\square$}}
\newcommand{\ms@minorstart}{%
  \pushQED{\qed}\renewcommand{\qedsymbol}{$\blacksquare$}}
\newenvironment{theorem}[1][]{\ms@majorstart\ms@theorem[#1]}
  {\popQED\endms@theorem}
\newenvironment{theorem*}[1][]{\ms@majorstart\ms@theoremstar[#1]}
  {\popQED\endms@theoremstar}
\newenvironment{proposition}[1][]{\ms@minorstart\ms@proposition[#1]}
  {\popQED\endms@proposition}
\newenvironment{lemma}[1][]{\ms@minorstart\ms@lemma[#1]}
  {\popQED\endms@lemma}
\newenvironment{definition}[1][]{\ms@minorstart\ms@definition[#1]}
  {\popQED\endms@definition}
\newenvironment{corollary}[1][]{\ms@majorstart\ms@corollary[#1]}
  {\popQED\endms@corollary}
\newenvironment{corollary*}[1][]{\ms@majorstart\ms@corollarystar[#1]}
  {\popQED\endms@corollarystar}
% The source's lemma/proposition variants are deliberately not carried over:
% choose one marker for each statement type rather than encode proof status.

\theoremstyle{remark}
\newtheorem{remark}[ms@proposition]{Remark}
\newtheorem*{remark*}{Remark}
% The ordinary amsthm proof environment retains its open square.

\crefname{ms@theorem}{theorem}{theorems}
\Crefname{ms@theorem}{Theorem}{Theorems}
\crefname{ms@proposition}{proposition}{propositions}
\Crefname{ms@proposition}{Proposition}{Propositions}
\crefname{ms@lemma}{lemma}{lemmas}
\Crefname{ms@lemma}{Lemma}{Lemmas}
\crefname{ms@definition}{definition}{definitions}
\Crefname{ms@definition}{Definition}{Definitions}
\crefname{ms@corollary}{corollary}{corollaries}
\Crefname{ms@corollary}{Corollary}{Corollaries}
\crefname{remark}{remark}{remarks}
\Crefname{remark}{Remark}{Remarks}

% Bibliography configuration for this manuscript.
  \RequirePackage[backend=biber,style=alphabetic,backref=true,
    giveninits=true,isbn=false,date=year]{biblatex}
  \addbibresource{\jobname.bib}
  \setcounter{biburlnumpenalty}{100}
  \setcounter{biburllcpenalty}{100}
  \setcounter{biburlucpenalty}{100}
  \renewcommand*{\bibfont}{\footnotesize}
  \renewbibmacro{in:}{}
  \DeclareFieldFormat{pages}{#1}
  % Use the public alphabetic style. No private author-hash overrides;
  % repeated author names are printed in full.
\makeatother
% END INLINE manuscript-style
\usepackage{mathrsfs}
% RNM1: paper-specific notation.
\newcommand{\R}{\mathbb R}
\newcommand{\C}{\mathbb C}
\newcommand{\Z}{\mathbb Z}
\newcommand{\N}{\mathbb N}
\newcommand{\dd}{\mathrm d}
% RNM1 / d038. Full proof companion to the RNO1 focused outline.
\title[Extraction and contraction of local interactions]
{Relevant extraction and contraction of local interactions}
\date{}
\begin{document}
\begin{abstract}
We study local Grassmann interactions on the three-dimensional
chamber lattice, with independent ordinary and covariant-difference
generators. Under explicit small-field and charge bounds, vacuum
and scalar mass extraction followed by a shortest-chain adjustment
makes the full ordinary quadratic kernel vanish at zero connection
on small polymers.
Scaling and reblocking give an explicit norm bound that is
contractive for sufficiently large block scale and sufficiently
small charge.
The low-degree gain follows from local gauge representatives with
an analytic disk of radius \(L^{7/4}\). We also prove a local
reconstruction estimate, uniform in the lattice mesh, using
cancellation in the averaged chamber Dirac operator.
\end{abstract}
\maketitle
\tableofcontents
\section{Introduction}
\label{sec:introduction}

Let \(L\ge5\) be odd, let \(M=L^m\) with \(m\in\Z_{\ge2}\), and put
\(\eta_k=L^{-k}\), \(k\in\Z_{\ge0}\). We use the photon torus
\(\eta_k\Z^3/\ell\Z^3\) and its chamber sites \((c,g)\), where \(g\)
is one of the \(48\) signed permutation matrices. Spinors have four
components and are antiperiodic in each coordinate. A polymer is a
face-connected union of the partition \(M\)-cubes. Its size parameter
\(Md_M(X)\) is the infimum of the lengths of connected finite nets
meeting every cube of \(X\); it is called small if \(d_M(X)\le L\).
The net convention and its estimates are given in
Appendix~\ref{app:polymer-geometry}.

An activity is a family of Grassmann polynomials in independent
ordinary generators \(\psi,\bar\psi\) and ordered-pair generators
\(\chi,\bar\chi\). The latter represent covariant H\"older differences
only when the fields are evaluated. We work with the space
\(K_k^{x,\C}(\ell)\) defined in
Section~\ref{sec:activities}: its families are balanced, holomorphic
on \(\widetilde R_k\times\mathfrak X_k\), local in the stated
ordinary and pair indices and in the connection bonds, and covariant
under the specified lattice, independent real-gauge and charge
actions. The real subspace \(K_k^x\) satisfies the additional
antilinear identity there. The two generator weights and family norm are
\[
 h_{k,1}=e_k^{-1/4},\qquad h_{k,2}=e_k^{-1/4+\varepsilon},\qquad
 \|\mathscr E\|_k^x
 =\sup_X e^{\kappa d_M(X)}
   \sup_{\widetilde R_k\times\mathfrak X_k}
                 \|\mathscr E(X)\|_{h_k}.
\]
Here \(0<e_k<1\); the weighted coefficient norm and the two
connection domains are specified in
\eqref{eq:rg-activity-coefficient-norm} and
\eqref{eq:connection-domains}. There is no smallness assumption
on the activity in the following linear theorem.

Set
\[
 \bar\varepsilon=\frac1{200},\quad c_X=54,\quad
 \kappa_0=54\log72,\quad K_{\mathrm{pol}}=\frac{e^{\kappa_0}}{36},\quad
 c_{\mathrm{rb}}=c_XK_{\mathrm{pol}},\quad c_{\mathrm{fl}}=61c_{\mathrm{rb}}.
\]
The scalar $K_{\mathrm{pol}}$ is the polymer-sum constant in
\eqref{eq:polymer-counting}.

\begin{theorem}[Extraction and normalized contraction]\label{thm:main}
Assume that both \(\ell/(LM)\) and \(\ell/(L^2M)\) are odd
integers at least \(13\). Suppose
\begin{equation}\label{eq:rg-contraction-parameters}
 \begin{gathered}
 0<\varepsilon\le\bar\varepsilon,\qquad 1-\varepsilon<\alpha<1,
 \qquad L^{1/8-3\bar\varepsilon/2}\ge2400,\qquad e_k^\varepsilon\le1/4,\\
 160L^{1+3\varepsilon/2}Me_k^\varepsilon\le1,\qquad
 9ML^\varepsilon e_k^{2\varepsilon}\le1,\qquad
 90c_XLMe_k^\varepsilon\le1,
 \end{gathered}
\end{equation}
and
\begin{equation}\label{eq:rg-contraction-polymer-parameters}
 \kappa\ge7\kappa_0+7,\qquad
 L(\kappa-6\kappa_0-4)-1\ge\kappa.
\end{equation}
Let \(\mathscr E\in K_k^{x,\C}(\ell)\) and
\(e_{k+1}=L^{1/2}e_k\).

\begin{enumerate}[label=\textup{(\roman*)},leftmargin=*]
\item There are complex-linear vacuum and mass maps, real on
\(K_k^x\), with
\[
 |\varepsilon_{\mathrm v}(\mathscr E)|
       \le K_{\mathrm{pol}}M^{-3}\|\mathscr E\|_k^x,\qquad
 |m(\mathscr E)|
       \le K_{\mathrm{pol}}M^{-3}e_k^{1/2}\|\mathscr E\|_k^x .
\]
The extracted mass matrix is \(m(\mathscr E)1_4\).
The local subtraction \(R\) of \eqref{eq:rg-normalized-remainder}
satisfies \(\|R\mathscr E\|_k^x\le2\|\mathscr E\|_k^x\) and the
exact independent-generator identity
\[
 \sum_X\mathscr E(X)
 =-\varepsilon_{\mathrm v}(\mathscr E)\ell^3
       -m(\mathscr E)\langle\bar\psi,\psi\rangle
       +\sum_X(R\mathscr E)(X).
\]
On a small polymer its zero-field vacuum kernel and the integral
of its zero-field ordinary quadratic kernel vanish.

\item The shortest-chain adjustment of
\eqref{eq:rg-adjustment-definition} gives
\(E=(R\mathscr E)^{\mathrm{nat}}\in K_k^{x,\C}(\ell)\), with
\(\|E\|_k^x\le20\|\mathscr E\|_k^x\). It leaves the evaluated
value of \(R\mathscr E\) unchanged whenever
\(\chi=\delta_{\alpha,\mathscr A+\mathscr C}\psi\) and
\(\bar\chi=\delta_{\alpha,-(\mathscr A+\mathscr C)}\bar\psi\).
On every small polymer it has the stronger normalization
\[
 E_{00}(X;0,0)=0,\qquad E_{20}(X;0,0;\cdot,\cdot)=0
\]
as full antisymmetric kernels. Both subtraction and adjustment
leave large-polymer activities unchanged.

\item Scale connections by
\(A_L(b)=L^{-1/2}A(b/L)\), ordinary generators by \(L^{-1}\),
and pair generators by \(L^{-1-\alpha}\).
Let \(\bar X^L\) be the union of \(LM\)-cubes containing \(X\), and set
\[
 (\mathcal LE)(Y;\mathscr A,\mathscr C)
       =\sum_{\bar X^L=LY}S_LE(X;\mathscr A_L,\mathscr C_L).
\]
Then \(\mathscr E\mapsto\mathcal L(R\mathscr E)^{\mathrm{nat}}\)
is a linear map into \(K_{k+1}^{x,\C}(\ell/L)\) and
\[
 \|\mathcal LE\|_{k+1}^x
       \le c_{\mathrm{fl}}L^{-1/4+3\varepsilon/2}\|\mathscr E\|_k^x.
\]
For \(w_j^{\mathrm E}=e_j^{1/4-7\varepsilon}\), the relative estimate is
\[
 \frac{\|\mathcal LE\|_{k+1}^x}{w_{k+1}^{\mathrm E}}
 \le c_{\mathrm{fl}}L^{-3/8+11\varepsilon/2}
                 \frac{\|\mathscr E\|_k^x}{w_k^{\mathrm E}}.
\]
Its factor is at most \(1/16\) if, in addition,
\(c_{\mathrm{fl}}L^{-3/8+6\bar\varepsilon}\le1/16\).
\end{enumerate}
All the stated maps preserve the real subspaces and commute with
setting the second connection argument to zero.
\end{theorem}

The proof reproduces the extraction, adjustment and contraction
arguments of \cite[\S12]{QED3Manuscript}. The first obstruction is
the ordinary quadratic term. Subtracting a local mass makes its
integral vanish, but does not make its kernel vanish at each pair
of points. We transport both legs along shortest chains of cubes
and average their terminal points. The remaining ordinary
quadratic term then factors through the integral of its kernel.
The other terms contain difference generators, whose scaling is
stronger. This construction also explains why the independent
generators are retained until evaluation.

The two lowest degrees require more than degree scaling. On a
small polymer, a real quadratic gauge removes the real constant
and symmetric linear parts of a scaled connection. Two cutoffs,
one for the Taylor remainder and one for the residual linear
field, produce global representatives in an analytic disk of
radius \(L^{7/4}\). The constant complex fields left over are
removed by holomorphic continuation of real-gauge identities.
The zero vacuum value and charge evenness give a second-order
Cauchy gain; the full quadratic normalization gives a first-order
gain. After reblocking, these estimates yield Theorem~\ref{thm:main}.
Sections~\ref{sec:extraction}--\ref{subsec:rg-activity-contraction}
give the proof.

Local reconstruction is a separate estimate, used when unit-scale
mean-field generators are replaced by ordinary and boundary-pair
legs. Theorem~\ref{thm:reconstruction}, stated with its own
hypotheses in Section~\ref{subsec:rg-field-reconstruction}, gives
uniform row bounds for that substitution. The apparent
\(\eta_k^{-1}\) loss in the Dirac stencil cancels before estimation.
The remaining phase defect is summed according to the first scale
at which two neighboring sites have the same block center. Recovery
of a mean field additionally assumes the fine operator is invertible
at the connection in question. No inverse is needed for the row
estimates or for Theorem~\ref{thm:main}.
\section{Activities, connections and coefficient norms}
\label{sec:activities}

We specify the local activity space used in Theorem~\ref{thm:main}.
The conventions are those of \cite[\S12.1]{QED3Manuscript}.
Only the finite Grassmann algebra, connection domains and symmetry
actions enter the extraction and contraction arguments.

\subsection{Chamber sites, pairs and connection domains}

Write \(O_{\mathrm{cub}}\) for the signed permutation matrices on
\(\R^3\). A chamber site is \(x=(c,g)\), with photon projection
\(\pi x=c\) and \(g\in O_{\mathrm{cub}}\). Integration over chamber
sites uses \(\dd x=(\eta_k^3/48)\sum_{c,g}\). We take
\(M\)-cubes centered at \(M\Z^3/\ell\Z^3\); their lattice sites
form a partition. Blocks of every odd side used below are centered
in the same way.

For the three internal neighbors let \(s_0\) change the sign of
the third coordinate, let \(s_1\) exchange the second and third
coordinates, and let \(s_2\) exchange the first and second.
Let \(r_{e_1}\) change the sign of the first coordinate. Set
\[
 x^{(i)}=(c,gs_i)\quad(0\le i\le2),\qquad
 x^{(3)}=(c+\eta_kge_1,gr_{e_1}).
\]
Each neighbor map is an involution. An oriented bond field
\(B\) is antisymmetric under reversal and periodic on the photon
lattice. Write \(B_\mu(c)=B(c,c+\eta_ke_\mu)\), and use the
forward differences
\[
 \partial_\nu B_\mu(c)=
       \eta_k^{-1}\bigl(B_\mu(c+\eta_ke_\nu)-B_\mu(c)\bigr),
 \qquad
 (d\lambda)_\mu(c)=
       \eta_k^{-1}\bigl(\lambda(c+\eta_ke_\mu)-\lambda(c)\bigr).
\]
Thus \((dB)_{\nu\mu}=\partial_\nu B_\mu-\partial_\mu B_\nu\).
Distances between photon sites are sup distances, using the
shortest lift when a distance bound is imposed. For \(0<\alpha<1\)
the three norms are
\[
 \|B\|_1=\|B\|_\infty,\qquad
 \|B\|_2=\|\partial B\|_\infty,\qquad
 \|B\|_3=\|\delta_\alpha\partial B\|_\infty
 =\sup_{\substack{\mu,\nu,c,c'\\0<d(c,c')\le1}}
 \frac{|\partial_\nu B_\mu(c')-\partial_\nu B_\mu(c)|}
      {d(c,c')^\alpha}.
\]
Local versions use only the indicated cells in the given region.
Let \(\square^\dagger\) be the concentric enlargement of an
\(M\)-cube to side \((2L+1)M\). Put
\(\theta_i=e_k^{-3/4+i\varepsilon}\).
The connection domains are
\begin{equation}\label{eq:connection-domains}
\begin{split}
 R_k&=\{B:\|B\|_i<\theta_i,\ 1\le i\le3\},\\
 \widetilde R_k&=\{B_0+B_1:B_1\in R_k,\ B_0\text{ real, and on each }
          \square^\dagger,\\
 &\hspace{25mm} B_0=B_\square^{\mathrm b}+d\lambda_\square
   \text{ with }B_\square^{\mathrm b},\lambda_\square\text{ real and }
          \|B_\square^{\mathrm b}\|_i<\theta_i\},\\
 \mathfrak X_k&=\{C:\|C\|_\infty<e_k^{-1+2\varepsilon},\
                       \|\partial C\|_\infty<e_k^{-1+3\varepsilon}\}.
\end{split}
\end{equation}
Only existence of these representatives is required.
For \(0<c\le1\), the notation \(c\widetilde R_k\) multiplies
all three bounds for both summands by \(c\).
For \(\widetilde R_k^-\) the local real bounds remain
\(\theta_i\), while the global summand has bounds
\((1-e_k^\varepsilon)\theta_i\).

For a short ordered pair \(x=(c,g),x'=(c',g')\), let
\(\tau B(c,c')\) be the arithmetic mean of the sums of \(B\)
along the six coordinate-order paths from \(c\) to \(c'\).
The bond values in this sum are unscaled; define
\[
 U_B(x,x')=\exp\{ie_k\eta_k\tau B(c,c')\},\qquad
 (\delta_{\alpha,B}\psi)(x,x')
       =\frac{U_B(x,x')\psi(x')-\psi(x)}
                    {\max(|c-c'|_\infty,\eta_k)^\alpha}.
\]
For pairs within one cluster the transport is one. Barred fields
use \(U_B^{-1}\). These formulas specify evaluation; the pair
generators in the algebra remain independent.

We use Pauli matrices and
\[
 \gamma_\mu=\sigma_3\otimes\sigma_\mu,\quad
 \gamma_5=\sigma_1\otimes1_2,\quad \Sigma_3=\sigma_3\otimes1_2,
 \qquad
 \mathsf C=i\gamma_2,\quad
 \varkappa=\gamma_5\mathsf C=\sigma_2\otimes\sigma_2 .
\]
Here \(\mathsf C\) is real with
\(\mathsf C^T=\mathsf C^{-1}=-\mathsf C\), while
\(\varkappa\) is a real symmetric involution. The spin group below
is the group generated by \(\gamma_1,\gamma_2,\gamma_3,\gamma_5\).
Its diagonal coordinate action \(r(S)\) is determined by
\(S\gamma_\mu S^{-1}=r(S)_{\mu\mu}\gamma_\mu\).
Connections transform by pullback on oriented bonds.

% QEM1:10142--10371; RG balls and mean-field specialization omitted.
\subsection{The local Grassmann algebra}
\label{subsec:rg-activity-spaces}

Fix a level $k$, mesh $\eta_k=L^{-k}$ and an admissible side
$\ell$, with $\ell/(LM)$ odd and at least $13$.
Let $\mathcal P_k$ be the collection of face-connected polymers
of $M$-cubes on this torus. A polymer is small when
$d_M(X)\le L$. A small polymer has an injective lift, unique up
to a period translation, by the diameter and lifting estimates
of Appendix~\ref{app:polymer-geometry}.
Write $\operatorname{Vol}(X)=M^3|X|_M$.

The independent Grassmann generators are $\psi$, $\bar\psi$ on
the chamber sites $x=(c,g)$ and $\chi$, $\bar\chi$ on ordered
pairs. For distinct chamber sites on the covering lattice, put
$d_X(x,x')=\max(|c-c'|_\infty,\eta_k)$ and restrict pairs to
$d_X\le M$. Pairs on the torus are simultaneous period classes;
their shortest lifts are unique since $2M<\ell$. A pair is
admissible for $X$ when its photon sites and all six rectilinear
paths between them lie in $X$, or when it is $(x,x^{(3)})$ and
one of its photon sites belongs to $X$.
Thus a pair may leave $X$ only along a bond of type $3$.
Pairs inside a cluster have $d_X=\eta_k$ and are admissible when
their common photon site belongs to $X$.
Ordinary generators acquire the sign $(-1)^{n_1+n_2+n_3}$
under translation by $\ell n$. Pair generators acquire the same
sign when both endpoints are translated by $\ell n$.
The kernels are antiperiodic in each ordinary argument and in
each simultaneous pair argument with these conventions.

Let $\xi=(x,\beta,\omega)$ index an ordinary leg and
$\zeta=(x,x',\beta,\omega)$ a H\"older leg, where
$\omega=0$ is unbarred and $\omega=1$ is barred.
We use the index measures
\[
 \nu(\xi)=\eta_k^3/48,\qquad
 \nu(\zeta)=(\eta_k^3/48)^2.
\]
Integrals sum the spin and bar indices as well. An activity
has a unique expansion with separately antisymmetric kernels,
\begin{equation}\label{eq:rg-activity-kernels}
 \mathscr E(X;\mathscr A,\mathscr C)
 =\sum_{n,m}\frac1{n!m!}\int\dd\xi\,\dd\zeta\,
      \mathscr E_{nm}(X;\mathscr A,\mathscr C;\xi;\zeta)
                       \psi(\xi_1)\cdots\psi(\xi_n)
                       \chi(\zeta_1)\cdots\chi(\zeta_m).
\end{equation}
Barred generators are included in this notation.
The generator weights, as distinct from the index measures, are
\[
 h_{k,1}=e_k^{-1/4},\qquad h_{k,2}=e_k^{-1/4+\varepsilon}.
\]
For an element $F$ the weighted coefficient norm is
\begin{equation}\label{eq:rg-activity-coefficient-norm}
 \|F\|_{h_k}
       =\sum_{n,m}\frac{h_{k,1}^nh_{k,2}^m}{n!m!}
                                   \int\dd\xi\,\dd\zeta\,|F_{nm}|.
\end{equation}
In particular a generator has norm $h_{k,i}$;
the weight $\nu$ enters only when coefficients are written as
integral kernels.

An activity family in $K_k^{x,\C}$ is balanced: every monomial
has as many barred as unbarred legs. Its kernels are holomorphic
on $\widetilde R_k\times\mathfrak X_k$, vanish unless all ordinary
legs lie over $X$ and all pairs are admissible for $X$, and read
the two connections only on bonds meeting $X$. The norm is
\begin{equation}\label{eq:rg-activity-norm}
 \|\mathscr E\|_k^x
  =\sup_{X\in\mathcal P_k}\exp(\kappa d_M(X))
    \sup_{(\mathscr A,\mathscr C)\in\widetilde R_k\times\mathfrak X_k}
                          \|\mathscr E(X;\mathscr A,\mathscr C)\|_{h_k}.
\end{equation}
We require $\kappa\ge\kappa_0$ and finite norm.

The required symmetries act on both the connections and the
generators. For $\tilde g=(S,a)$ with the coordinate sign action
$r(S)$ and $a\in M\Z^3$, put
\[
 \tilde g(c,g)=(r(S)c+a,r(S)g),\qquad
 a_{\tilde g}\psi(x)=S\psi(\tilde g^{-1}x),\qquad
 a_{\tilde g}\bar\psi(x)=(S^{-1})^T\bar\psi(\tilde g^{-1}x).
\]
The same spin matrices act on pair generators after transforming
both endpoints. Local gauge transformations multiply an unbarred
ordinary or pair leg by $\exp(ie_k\Lambda(c))$, using the photon
site of its first point, and a barred leg by the inverse phase.
Write this automorphism as $G_\Lambda$.
Charge conjugation and the antilinear reality map act by
\[
 \mathrm{Cg}\,\psi=\mathsf C\bar\psi,\qquad
 \mathrm{Cg}\,\bar\psi=-\mathsf C\psi,\qquad
 \mathbf K\psi=\varkappa\psi,\qquad
 \mathbf K\bar\psi=\varkappa\bar\psi,
\]
and by the same rules on the H\"older generators.

The covariance requirements on a family are
\begin{equation}\label{eq:rg-activity-symmetries}
 \begin{split}
 a_{\tilde g}\mathscr E(\tilde gX;\tilde g\mathscr A,
                                      \tilde g\mathscr C)
                                  &=\mathscr E(X;\mathscr A,\mathscr C),\\
 \mathrm{Cg}\,\mathscr E(X;-\mathscr A,-\mathscr C)
                                  &=\mathscr E(X;\mathscr A,\mathscr C).
 \end{split}
\end{equation}
For real scalar functions $\lambda,\lambda'$ on the photon torus
and two connection pairs in the domain
which satisfy $\mathscr A^1=\mathscr A-d\lambda$ and
$\mathscr C^1=\mathscr C-d\lambda'$ on bonds at distance at
most one from $X$, the local gauge requirement is
\[
 G_{\lambda+\lambda'}\mathscr E(X;\mathscr A^1,\mathscr C^1)
                               =\mathscr E(X;\mathscr A,\mathscr C).
\]
The real subspace $K_k^x$ has in addition
\[
 \mathbf K\mathscr E(X;-\overline{\mathscr A},-\overline{\mathscr C})
                               =\mathscr E(X;\mathscr A,\mathscr C).
\]
All requirements refer to the generators before evaluating fields.
They are essential to the form of the counterterms.

The spaces $K_k^\C,K_k$ without the second connection argument
have the same definitions on $\widetilde R_k$. Restriction
$\mathscr E|_0(X;\mathscr A)=\mathscr E(X;\mathscr A,0)$ maps
$K_k^{x,\C}$ into $K_k^\C$ and has norm at most one. The real
subspaces also restrict. Analyticity, locality and the symmetries
follow by setting the second argument and its gauge function to
zero; the norm bound follows from the smaller supremum.

\subsection{Total connections}
\label{subsec:rg-total-connection-domain}

The comparison domain $\widetilde G_k$ consists of fields $B_0+B_1$
with $B_0$ real and $\|B_1\|_\infty<e_k^{-1+\varepsilon}$,
where on each $\square^{\sim5}$ the field $B_0$ has a real-gauge
representative of supremum norm below $e_k^{-1+\varepsilon}$.
Here $\square^{\sim5}$ has the same center and side $11M$.
The real-gauge representatives in the domains defined above may
be chosen independently on each enlarged cube.

Assume $L\ge5$, $0<\varepsilon\le1/200$ and
\begin{equation}\label{eq:rg-domain-charge-conditions}
 e_k^\varepsilon\le\tfrac14,\qquad
 e_k^{1/4}+e_k^\varepsilon\le1,\qquad
 4e_k^{1/4+2\varepsilon}+2e_k^{3\varepsilon}\le1.
\end{equation}
Then
\begin{equation}\label{eq:rg-total-domain-inclusion}
 \widetilde R_k+\mathfrak X_k\subset\widetilde G_k.
\end{equation}
Indeed, write $\mathscr A=A_0+A_1$ with $A_1\in R_k$ and
$A_0$ locally real-gauge equivalent on $\square^\dagger$ to
$A_\square^b$ with the three threshold bounds. Restrict the
representative to $\square^{\sim5}\subset\square^\dagger$.
Its supremum norm is below
$e_k^{-3/4+\varepsilon}\le e_k^{-1+\varepsilon}$, and
\[
 \|A_1+\mathscr C\|_\infty
 <e_k^{-3/4+\varepsilon}+e_k^{-1+2\varepsilon}
 =(e_k^{1/4}+e_k^\varepsilon)e_k^{-1+\varepsilon}
 \le e_k^{-1+\varepsilon}.
\]
This is the required decomposition of the total connection.
Real gauge transformations can absorb any real constant part,
locally on the nonwrapping cubes.

Put $\mathfrak i(B)=e_k\|\operatorname{Im}B\|_\infty$ and
$q_d(B)=e_k\|dB\|_\infty$. The same decomposition gives
\begin{equation}\label{eq:rg-total-imaginary-curvature-bounds}
 \mathfrak i(\mathscr A+\mathscr C)
       \le e_k^{1/4+\varepsilon}+e_k^{2\varepsilon}\le1,
 \qquad
 q_d(\mathscr A+\mathscr C)
       \le4e_k^{1/4+2\varepsilon}+2e_k^{3\varepsilon}\le1.
\end{equation}
For the curvature, every plaquette lies in an enlarged cube,
and $dA_0=dA_\square^b$ there. This costs at most
$2\|\partial A_\square^b\|_\infty$; the global small field
costs another $2\|\partial A_1\|_\infty$, and the external
field costs $2\|\partial\mathscr C\|_\infty$.
The imaginary bound uses only that $A_0$ is real.

The domains are open: the strict inequalities define open
sets for the global summand, which can absorb a sufficiently
small complex perturbation while the local real representative
is kept fixed. The norm-defined sets are invariant under sign
and conjugation; the local real-gauge definitions are also
invariant under real gauges and real constant additions.
Coordinate sign changes and $M$ translations preserve them.
Their derivative and H\"older norms are preserved, with a
possible reversal of a bond and translation of its base point.
Thus \eqref{eq:rg-total-domain-inclusion} applies also at the
signs, conjugates and real gauges used below.

\subsection{Weighted linear substitution}

We will use the following finite-algebra fact, namely the
linear-substitution part of
\cite[Proposition~11.4]{QED3Manuscript}. Positive index measures
do not change it: when a polynomial is written in ordered
monomials, each coefficient already includes the product of its
index measures.

\begin{lemma}\label{lem:linear-substitution}
Give each generator of a finite Grassmann algebra a positive
weight, and give the algebra the sum of absolute ordered
coefficients times the corresponding products of weights.
This norm is submultiplicative. A homomorphism that sends each
generator to an odd linear combination of target generators
with norm at most its original weight decreases this norm.
\end{lemma}

\begin{proof}
In multiplying two ordered monomials, a repeated generator makes
the product zero. Otherwise the product differs from the ordered
union only by a sign, and its weight is the product of the two
weights. The triangle inequality, summed over the two sets of
coefficients, proves submultiplicativity. The norm of the image
of each ordered monomial is consequently at most its original
weight. Summing its absolute coefficient proves the
substitution bound. Separately antisymmetric kernels give the
same coefficient sum: the factorials cancel the number of
orders of each set of distinct indices.
\end{proof}
\section{Extraction and quadratic normalization}
\label{sec:extraction}

The extraction and shortest-chain adjustment follow
\cite[\S\S12.3,12.5]{QED3Manuscript}. They solve different problems.
The extracted matrix becomes scalar only after summing over polymers.
The adjustment then converts the vanishing integral of each remaining
quadratic kernel into vanishing of the full kernel at zero connection.
This is the normalization used in the Cauchy estimate.

\subsection{Vacuum and mass extraction}
\label{subsec:rg-extraction}

For a balanced activity, its part of bidegree $(2,0)$ can be
written uniquely as
\[
 \mathscr E^{[2,0]}(X;\mathscr A,\mathscr C)
   =\int\dd x\,\dd x'\,
             \bar\psi(x)^T K_X(\mathscr A,\mathscr C;x,x')\psi(x').
\]
The matrix kernel is the ordinary antisymmetric kernel with
the first bar index one and the second bar index zero.
Antisymmetry accounts for the other ordering; in particular
\begin{equation}\label{eq:rg-bilinear-norm}
 \|\mathscr E^{[2,0]}(X;\mathscr A,\mathscr C)\|_{h_k}
   =h_{k,1}^2\sum_{\beta,\beta'}\int\dd x\,\dd x'
                     |K_X(\mathscr A,\mathscr C;x,x')_{\beta\beta'}|.
\end{equation}
The kernel is antiperiodic in each argument separately.

For a small $X$, choose one lift $\widetilde X$ and define
\begin{equation}\label{eq:rg-local-extractions}
 a_0(\mathscr E,X)
       =\operatorname{Vol}(X)^{-1}\mathscr E_{00}(X;0,0),\qquad
 a_2(\mathscr E,X)
       =\operatorname{Vol}(X)^{-1}
             \int_{\widetilde X}\dd x\int_{\widetilde X}\dd x'\,
                                                    K_X(0,0;x,x').
\end{equation}
Both cluster labels are averaged with measure $\eta_k^3/48$.
The second expression is independent of the lift: simultaneous
translation by a period multiplies its two arguments by the
same antiperiodic sign, whose product is one.

For an $M$-cube $\square$ put
\begin{equation}\label{eq:rg-counterterm-maps}
 \varepsilon_{\mathrm v}(\mathscr E,\square)
          =-\sum_{\substack{X\supset\square\\X\text{ small}}}
                                                     a_0(\mathscr E,X),
 \qquad
 M(\mathscr E,\square)
          =-\sum_{\substack{X\supset\square\\X\text{ small}}}
                                                     a_2(\mathscr E,X).
\end{equation}
Translations by $M\Z^3$ identify the small-polymer sets over
different cubes. They leave the two extracted quantities
unchanged. Thus the two sums do not depend on $\square$.
Write their values as $\varepsilon_{\mathrm v}(\mathscr E)$ and
$M(\mathscr E)$.

\begin{proposition}\label{prop:rg-scalar-counterterms}
There is a complex-linear scalar map $m(\mathscr E)$ such that
$M(\mathscr E)=m(\mathscr E)1$. Both $m$ and
$\varepsilon_{\mathrm v}$ are real on $K_k^x$, and
\begin{equation}\label{eq:rg-counterterm-bounds}
 |\varepsilon_{\mathrm v}(\mathscr E)|
                  \le K_{\mathrm{pol}}M^{-3}\|\mathscr E\|_k^x,
 \qquad
 |m(\mathscr E)|\le K_{\mathrm{pol}}M^{-3}e_k^{1/2}\|\mathscr E\|_k^x.
\end{equation}
They commute with restriction of the second connection argument
to zero.
\end{proposition}

\begin{proof}
Substitution of the generator transformations in a bilinear,
followed by uniqueness of its kernel, gives
\begin{equation}\label{eq:rg-extraction-matrix-symmetries}
 \begin{split}
 a_2(\mathscr E,\tilde gX)&=S a_2(\mathscr E,X)S^{-1},\\
 a_2(\mathscr E,X)&=\mathsf C^{-1}a_2(\mathscr E,X)^T\mathsf C,\\
 a_2(\mathscr E,X)&=\varkappa^{-1}
                         \overline{a_2(\mathscr E,X)}\varkappa
                                                   \quad(\mathscr E\in K_k^x).
 \end{split}
\end{equation}
For the first relation, the two variables and their labels
are transformed together, and the measure is unchanged.
For charge conjugation, moving the unbarred leg past the barred
leg cancels the minus sign in its transformation, and exchanges
the two kernel variables. For the last relation, the reality
map conjugates the coefficient and multiplies both legs by
$\varkappa$; its inverse transpose is $\varkappa^{-1}$.
The constant kernel has the corresponding scalar invariance
and is real for a real family.

Let $z$ be the center of $\square$ and take
$S=\gamma_1,\gamma_2,\gamma_3$ or $\gamma_5$.
The symmetry $(S,(1-r(S))z)$ fixes $\square$ and permutes the
small polymers containing it. The first relation in
\eqref{eq:rg-extraction-matrix-symmetries} therefore shows that
$M(\mathscr E)$ commutes with all four matrices. Their commutant
is scalar. In the fixed Dirac basis, $\gamma_3$ and
$\Sigma_3=-i\gamma_1\gamma_2\gamma_3$ are diagonal with the
four distinct pairs of diagonal entries
$(1,1),(-1,1),(-1,-1),(1,-1)$. A commuting matrix is diagonal;
commutation with $\gamma_1$ identifies its first two and last
two entries, and commutation with $\gamma_5$ identifies these
pairs. This proves the scalar assertion. The last relation in
\eqref{eq:rg-extraction-matrix-symmetries} makes this scalar
real for a real family.

The coefficient norm gives, for each small $X$,
\begin{equation}\label{eq:rg-single-extraction-bounds}
 \begin{split}
 \operatorname{Vol}(X)|a_0(\mathscr E,X)|
                    &\le\|\mathscr E(X;0,0)\|_{h_k},\\
 \operatorname{Vol}(X)\sum_{\beta,\beta'}
                              |a_2(\mathscr E,X)_{\beta\beta'}|
                    &\le h_{k,1}^{-2}
                               \|\mathscr E^{[2,0]}(X;0,0)\|_{h_k}.
 \end{split}
\end{equation}
Use $h_{k,1}^{-2}=e_k^{1/2}$ and
$\operatorname{Vol}(X)\ge M^3$. The individual vacuum and mass
contributions are bounded by $M^{-3}\|\mathscr E\|_k^x$ and
$M^{-3}e_k^{1/2}\|\mathscr E\|_k^x$, respectively, times
$\exp(-\kappa d_M(X))$.
The polymer sum over a fixed cube is at most $K_{\mathrm{pol}}$ by
\eqref{eq:polymer-counting} and $\kappa\ge\kappa_0$. For the mass use one diagonal entry of
the matrix sum, now known to be scalar. This proves the bounds.
Linearity follows from uniqueness of the antisymmetric kernels.
Restriction gives the same kernels at $(0,0)$, proving the
last assertion.
\end{proof}

Subtract these local terms only on small polymers:
\begin{equation}\label{eq:rg-normalized-remainder}
 (R\mathscr E)(X;\mathscr A,\mathscr C)
 =\mathscr E(X;\mathscr A,\mathscr C)
       -a_0(\mathscr E,X)\operatorname{Vol}(X)
       -\int_X\dd x\,\bar\psi(x)^Ta_2(\mathscr E,X)\psi(x).
\end{equation}
For large $X$ set $R\mathscr E=\mathscr E$. Then
\begin{equation}\label{eq:rg-exact-extraction-identity}
 \sum_X\mathscr E(X;\mathscr A,\mathscr C)
  =-\varepsilon_{\mathrm v}(\mathscr E)\ell^3
     -m(\mathscr E)\langle\bar\psi,\psi\rangle
                           +\sum_X(R\mathscr E)(X;\mathscr A,\mathscr C).
\end{equation}
Indeed, decompose the integral and volume of each small polymer
over its $M$-cubes, interchange the finite sums, and use
\eqref{eq:rg-counterterm-maps} on each cube. Every chamber
component receives the same scalar mass coefficient:
\[
 m(\mathscr E)\langle\bar\psi,\psi\rangle
 =m(\mathscr E)\frac{\eta_k^3}{48}
         \sum_c\sum_g\sum_{\beta=1}^4
                                      \bar\psi_\beta(c,g)\psi_\beta(c,g).
\]

The diagonal bilinear of a constant matrix $A$ on $X$ has
kernel $48\eta_k^{-3}1[x=x'\text{ over }X]A$ when both
arguments are taken over the same lift. Its integral is
$\operatorname{Vol}(X)A$, and its norm is
\[
 h_{k,1}^2\operatorname{Vol}(X)
                         \sum_{\beta,\beta'}|A_{\beta\beta'}|.
\]
These identities and \eqref{eq:rg-single-extraction-bounds}
show that
\begin{equation}\label{eq:rg-remainder-properties}
 \begin{split}
 (R\mathscr E)_{00}(X;0,0)&=0\quad(X\text{ small}),\\
 \int_{\widetilde X}\dd x\,\dd x'\,
                  (R\mathscr E)_{20}(X;0,0;x,x')&=0\quad(X\text{ small}),\\
 \|R\mathscr E\|_k^x&\le2\|\mathscr E\|_k^x.
 \end{split}
\end{equation}
For the bound, the two subtracted terms have distinct bidegrees
and together have norm at most $\|\mathscr E(X;0,0)\|_{h_k}$.
Apply the triangle inequality to
\eqref{eq:rg-normalized-remainder}, then the activity norm.

The map $R$ preserves the activity spaces and restriction to
zero. The subtracted kernels are independent of the connections,
and their two legs lie over $X$, so analyticity and locality
remain. A diagonal bilinear is invariant under every local
gauge. Its lattice, charge and reality transformations follow
from \eqref{eq:rg-extraction-matrix-symmetries}; the volume
and constant term have the same symmetries. This verifies
each requirement on the remainder, without assuming that
an individual matrix $a_2(\mathscr E,X)$ is scalar.

For later comparisons, the extraction has a precise local
meaning across different sides. Fix a lifted $M$-cube and
consider all lifted small polymers containing it. Their
projections are bijective onto the small polymers containing
its image, and their $d_M$ and volumes are unchanged.
If two families have the same constant and quadratic kernels
on every such lifted polymer at zero connections, then their
vacuum and mass matrices over this cube agree term by term in
\eqref{eq:rg-counterterm-maps}. Volume independence of a flow
will require this equality of kernels from the step comparison.
\subsection{Adjustment of the quadratic kernel}
\label{subsec:rg-adjustment}

The extraction makes the integral of the zero-field quadratic
kernel vanish. For contraction we need the kernel itself to
vanish. We achieve this by writing each ordinary leg as a
transported average and a sum of H\"older differences. This
changes the activity before evaluation and leaves its evaluated
value unchanged.

Let $X$ be small, choose a lift $\widetilde X$, and let
$x=(c,g)$ and $\tilde x=(\tilde c,\tilde g)$ lie over this lift.
Take all shortest face chains of its $M$-cubes between their
two cubes. For a chain $(\square_0,\ldots,\square_n)$ with
centers $z_i$, use the points
\[
 p_0=x,\qquad p_{i+1}=(z_i,g)\quad(0\le i\le n),\qquad
 p_{n+2}=\tilde x.
\]
The initial label is kept at each intermediate center.
Each distinct consecutive pair is admissible for $X$ of the
first kind. The paths in the first and last cube remain in
those cubes; between adjacent centers all six paths are the
straight segment across their common face. The last pair
can have different labels, and can lie inside one cluster.
For a coincident pair there is no difference term.

With $u_B(p,p')=U_B(p,p')$ and $u_B(p,p)=1$, set
\[
 U_i^\pm=\prod_{j<i}u_B(p_j,p_{j+1})^{\pm1}.
\]
The chain has no repeated cube, and
$n+1\le|X|_M\le c_X(L+1)$. Hence
\begin{equation}\label{eq:rg-chain-length}
 \sum_{i=0}^{n+1}|\pi p_i-\pi p_{i+1}|_1\le(n+3)M
          \le3c_XLM,\qquad
 \sum_{i:p_i\ne p_{i+1}}d_X(p_i,p_{i+1})^\alpha\le3c_XLM.
\end{equation}
The second estimate uses $d_X^\alpha\le\max(1,d_X)\le M$.
The transport bound is
$|u_B(p,p')^{\pm1}|\le\exp(e_k\|\operatorname{Im}B\|_\infty
|\pi p-\pi p'|_1)$, obtained by counting the bonds in each path.

For $B=\mathscr A+\mathscr C$, the domain decomposition and
$e_k^\varepsilon\le1/4$ give
$e_k\|\operatorname{Im}B\|_\infty\le e_k^\varepsilon/2$.
Assume in addition
\begin{equation}\label{eq:rg-adjustment-charge-condition}
 90c_XLM e_k^\varepsilon\le1.
\end{equation}
The modulus of every partial product $U_i^\pm$ is then at most
$\exp(1)$, uniformly in the chain and mesh.

Let $\operatorname{Mn}$ be the mean over the shortest chains.
Define the homomorphism $\Theta_{X,B}$ on ordinary generators by
\begin{equation}\label{eq:rg-chain-homomorphism}
 \begin{split}
 \Theta_{X,B}\psi(x)
  =\operatorname{Vol}(X)^{-1}\int_{\widetilde X}\dd\tilde x\,
    \operatorname{Mn}\left[
       U_{n+2}^+\psi(\tilde x)
       -\sum_{i:p_i\ne p_{i+1}}
           U_i^+d_X(p_i,p_{i+1})^\alpha\chi(p_i,p_{i+1})\right],\\
 \Theta_{X,B}\bar\psi(x)
  =\operatorname{Vol}(X)^{-1}\int_{\widetilde X}\dd\tilde x\,
    \operatorname{Mn}\left[
       U_{n+2}^-\bar\psi(\tilde x)
       -\sum_{i:p_i\ne p_{i+1}}
           U_i^-d_X(p_i,p_{i+1})^\alpha\bar\chi(p_i,p_{i+1})\right].
 \end{split}
\end{equation}
The spin components are unchanged. Both displayed images are
linear combinations of odd independent generators, so they
anticommute and determine a unique homomorphism on the
subalgebra of ordinary legs over $X$. Translation of the lift
by a period preserves every transport and chain and multiplies
each generator by the same antiperiodic sign. Thus the
homomorphism is independent of this choice of lift.

If the generators are evaluated at fields satisfying their
covariant H\"older relations, repeated substitution of
\[
 \psi(p_i)=u_B(p_i,p_{i+1})\psi(p_{i+1})
                    -d_X(p_i,p_{i+1})^\alpha\chi(p_i,p_{i+1})
\]
along a chain gives exactly the first expression in
\eqref{eq:rg-chain-homomorphism}; the barred relation uses
$u_B^{-1}$. Taking the two means therefore reproduces the
original leg. This proves equality after evaluation for
every polynomial in the ordinary legs over $X$.

The adjusted family is
\begin{equation}\label{eq:rg-adjustment-definition}
 \mathscr E^{\mathrm{nat}}(X)
   =\mathscr E(X)-\mathscr E^{[2,0]}(X)
           +\Theta_{X,\mathscr A+\mathscr C}
                                      \mathscr E^{[2,0]}(X)
                                      \qquad(X\text{ small}),
\end{equation}
and is unchanged on large polymers. The equality just proved
shows that the adjusted and original families have equal
evaluated values. Only bidegrees $(2,0),(1,1),(0,2)$ can change.

At zero connection all $U_i^\pm$ are one. The ordinary part
of an adjusted bilinear with kernel $K$ is therefore
\[
 \bar\psi_{\mathrm{av}}^T
    \left[\int_{\widetilde X}\dd x\,\dd x'\,K(x,x')\right]
                                                     \psi_{\mathrm{av}},
 \qquad
 \psi_{\mathrm{av}}=\operatorname{Vol}(X)^{-1}
                                        \int_{\widetilde X}\dd x\,\psi(x).
\]
Apply this to $R\mathscr E$, whose kernel integral vanishes.
With $E=(R\mathscr E)^{\mathrm{nat}}$, we obtain the stronger
normalization
\begin{equation}\label{eq:rg-adjusted-kernel-normalization}
 E_{00}(X;0,0)=0,\qquad E_{20}(X;0,0;\cdot,\cdot)=0
                                                     \quad(X\text{ small}).
\end{equation}

The same chain estimate controls the norm of adjustment.
Decompose $\Theta\psi=L_\psi+L_\chi$ by bidegree. For arbitrary
generator weights $(w_1,w_2)$,
\[
 \|L_\psi\|_w\le\exp(1)w_1,\qquad
 \|L_\chi\|_w\le\exp(1)(3c_XLM)w_2.
\]
The normalized site and label mean has total mass one.
Writing $s_K=\sum_{\beta,\beta'}\int|K_{\beta\beta'}|$ for
the balanced quadratic kernel and $F=\Theta\mathscr E^{[2,0]}$
gives the three bounds
\[
 \begin{aligned}
 \|F^{[2,0]}\|_w&\le\exp(2)w_1^2s_K,\\
 \|F^{[1,1]}\|_w&\le2\exp(2)(3c_XLM)w_1w_2s_K,\\
 \|F^{[0,2]}\|_w&\le\exp(2)(3c_XLM)^2w_2^2s_K.
 \end{aligned}
\]
Since $h_{k,2}=e_k^\varepsilon h_{k,1}$ and
$3c_XLM e_k^\varepsilon\le1/30$, summing these bounds yields
\begin{equation}\label{eq:rg-adjustment-norm}
 \|\mathscr E^{\mathrm{nat}}(X)\|_{h_k}
 \le\exp(2)(1+3c_XLM e_k^\varepsilon)^2\|\mathscr E(X)\|_{h_k}
                                               \le10\|\mathscr E(X)\|_{h_k}.
\end{equation}
Subtracting the original quadratic part before adding its
image uses additivity of the norm over bidegrees.

All new legs lie over $X$ or at its admissible chain pairs,
and all new transports read bonds with both ends in $X$.
Thus adjustment preserves locality. It preserves analyticity
because these transports are entire in the connection.
For the lattice covariance, symmetries carry shortest chains,
their centers and labels to the corresponding image chains.
For a real gauge the factors along a partial chain telescope
to the phase difference at its endpoints, giving the original
leg phase. Charge conjugation replaces $u_B$ by $u_B^{-1}$;
the reality map uses $\overline{u_{-\overline B}}=u_B$.
These identities on generators prove all four covariance
identities for $\Theta$, and hence the activity symmetries
for adjustment. Restriction to $\mathscr C=0$ commutes with
adjustment because it leaves the same total connection and
the same ordinary quadratic kernel. Therefore adjustment
preserves the complex and real activity spaces.
\section{Scaling and reblocking}\label{subsec:rg-scaling}

Throughout Sections~\ref{subsec:rg-scaling}--\ref{subsec:rg-activity-contraction},
both $\ell$ and $\ell/L$ are admissible and the hypotheses of
Theorem~\ref{thm:main} hold. In particular $\ell\ge13L^2M$.
We first record the exact scaling identities from
\cite[\S12.6]{QED3Manuscript}.

Write the three thresholds of $R_k$ as
\[
 \theta_1=e_k^{-3/4+\varepsilon},\qquad
 \theta_2=e_k^\varepsilon\theta_1,\qquad
 \theta_3=e_k^\varepsilon\theta_2,
\]
and the thresholds of $\mathfrak X_k$ as
$d_1=e_k^{-1+2\varepsilon}$ and $d_2=e_k^\varepsilon d_1$.
For a level $k+1$ bond function and gauge function define
\begin{equation}\label{eq:rg-connection-scaling}
 A_L(b)=L^{-1/2}A(b/L),\qquad
 \lambda_L(c)=L^{1/2}\lambda(c/L).
\end{equation}
Since $\eta_{k+1}=\eta_k/L$, the finite differences obey
$(d\lambda)_L=d\lambda_L$ and
\[
 \|A_L\|_\infty=L^{-1/2}\|A\|_\infty,\qquad
 \|\partial A_L\|_\infty=L^{-3/2}\|\partial A\|_\infty,\qquad
 \|\delta_\alpha\partial A_L\|_\infty
                \le L^{-3/2-\alpha}\|\delta_\alpha\partial A\|_\infty.
\]
The last inequality uses only the coarse pairs at distance
$\le1/L$, a subset of those in the coarse seminorm.
The charge relation $e_{k+1}=L^{1/2}e_k$ therefore gives
the following fine-scale bounds for the scaled complex
small field, and for each scaled local real representative:
\begin{equation}\label{eq:rg-scaled-field-thresholds}
 \begin{split}
 b_1&=L^{-7/8+\varepsilon/2}\theta_1,\\
 b_2&=L^{-15/8+\varepsilon}\theta_2,\\
 b_3&=L^{-15/8+3\varepsilon/2}\theta_3.
 \end{split}
\end{equation}
Here the additional factor $L^{-\alpha}$ in the third bound
has been discarded. For the external argument the bounds are
\begin{equation}\label{eq:rg-scaled-external-thresholds}
 \|\mathscr C_L\|_\infty<L^{-1+\varepsilon}d_1,\qquad
 \|\partial\mathscr C_L\|_\infty<L^{-2+3\varepsilon/2}d_2.
\end{equation}
The enlarged coarse cube used to define a local representative,
after multiplication by $L$, contains the enlarged fine cube.
Scaling its real gauge thus gives the corresponding fine
representative. All the ratios in
\eqref{eq:rg-scaled-field-thresholds} and
\eqref{eq:rg-scaled-external-thresholds} are at most $L^{-3/4}$,
by \eqref{eq:rg-contraction-parameters}. It follows that
\begin{equation}\label{eq:rg-scaled-domain-inclusion}
 \mathscr A_L\in L^{-3/4}\widetilde R_k
                     \subset\widetilde R_k^-,\qquad
 \mathscr C_L\in L^{-3/4}\mathfrak X_k.
\end{equation}
For the second inclusion in the first statement use
$L^{-3/4}\le1/2\le1-e_k^\varepsilon$.

The Grassmann scaling homomorphism is
\begin{equation}\label{eq:rg-leg-scaling}
 S_L\psi(x)=L^{-1}\psi'(x/L),\qquad
 S_L\chi(x,x')=L^{-1-\alpha}\chi'(x/L,x'/L),
\end{equation}
with the same powers for the barred generators.
Only the photon coordinate is divided by $L$; the chamber
label and spin index are kept. Ordinary indices are mapped
bijectively onto the coarse ordinary indices. Pair indices
are mapped bijectively onto the coarse pair indices with
$d_X\le M/L$, since $d_X$ is divided by $L$.
This includes a pair inside one cluster:
$d_X=\eta_k$ becomes $d_X=\eta_{k+1}$.
The antiperiodic signs agree under this bijection.

Thus distinct basis monomials have distinct images, up to
their ordering signs, and every monomial of bidegree $(n,m)$
is multiplied by $L^{-n-(1+\alpha)m}$. There is no cancellation
of coefficients. For arbitrary positive weights,
\begin{equation}\label{eq:rg-scaling-norm-identity}
 \|S_LF\|_{w_1,w_2}=\|F\|_{L^{-1}w_1,L^{-1-\alpha}w_2}.
\end{equation}
As $h_{k+1,1}/h_{k,1}=L^{-1/8}$ and
$h_{k+1,2}/h_{k,2}=L^{-1/8+\varepsilon/2}$, this gives
\begin{equation}\label{eq:rg-degree-scaling}
 \|(S_LF)^{[n,m]}\|_{h_{k+1}}
       =L^{-p(n,m)}\|F^{[n,m]}\|_{h_k},\qquad
 p(n,m)=n+(1+\alpha)m+\frac{n+m}{8}-\frac{m\varepsilon}{2}.
\end{equation}
The powers are identical for all admissible pair types.

In kernel notation the same calculation reads
\begin{equation}\label{eq:rg-scaled-kernel}
 (S_LF)_{nm}(\xi';\zeta')=
 \begin{cases}
 L^{3n+6m-n-(1+\alpha)m}F_{nm}(L\xi';L\zeta'),
                       &d_X(\zeta'_j)\le M/L\text{ for every }j,\\
 0,                    &\text{otherwise}.
 \end{cases}
\end{equation}
Indeed, the index weights satisfy
$\eta_k^3/48=L^3\eta_{k+1}^3/48$ and hence contribute
$L^{3n+6m}$ to the kernel. Antisymmetry and antiperiodicity
are preserved, so uniqueness of the kernel expansion gives
\eqref{eq:rg-scaled-kernel}.

Let $\bar X^L$ be the union of $LM$-cubes containing the cubes
of $X$. Define the reblocked family by
\begin{equation}\label{eq:rg-reblocking-definition}
 (\mathcal L F)(Y;\mathscr A,\mathscr C)
       =\sum_{X:\bar X^L=LY}S_L
                          F(X;\mathscr A_L,\mathscr C_L).
\end{equation}
The sum is finite on a fixed torus. It is local in $Y$ if
$F$ is local in each $X$. Ordinary legs over $X$ scale to
legs over $Y$; the six paths for a pair of the first kind
scale to paths in $Y$. For a boundary pair,
$x^{(3)}/L=(x/L)^{(3)}$, and an endpoint in $X$ scales to
an endpoint in $Y$. Finally, every fine bond meeting $X$
scales to a coarse bond meeting $Y$. Equality of both
coarse connection arguments on these bonds consequently
gives equality of the fine arguments read by each term.
This proves both the support and the coefficient locality
of \eqref{eq:rg-reblocking-definition}.
\section{Global representatives on a small polymer}
\label{subsec:rg-small-polymer-gauge}

The construction in this section follows \cite[\S12.7]{QED3Manuscript}.
The scaled fields belong to the fine domain, but this alone
does not give an analytic gain for their vacuum and ordinary
quadratic kernels. On a small polymer they admit a stronger
representation, with a real gauge, two constant fields and
global fields small enough for an analytic disk of radius
$L^{7/4}$.

\begin{proposition}\label{prop:rg-small-polymer-representatives}
Assume \eqref{eq:rg-contraction-parameters}. Let $X$ be small
and let $(\mathscr A,\mathscr C)\in
\widetilde R_{k+1}\times\mathfrak X_{k+1}$.
There are a real function $\lambda_*$, constant bond
functions $\mathscr K_1,\mathscr K_2$, and global bond
functions $\mathscr A^{\mathrm b},\mathscr C^{\mathrm b}$
on the fine torus such that, on every bond with an endpoint
at distance at most $1$ from $X$,
\begin{equation}\label{eq:rg-small-polymer-gauge-identity}
 \mathscr A_L-d\lambda_*=\mathscr A^{\mathrm b}+\mathscr K_1,
 \qquad
 \mathscr C_L=\mathscr C^{\mathrm b}+\mathscr K_2.
\end{equation}
They satisfy
\begin{equation}\label{eq:rg-small-polymer-global-bounds}
 \begin{split}
 \|\mathscr A^{\mathrm b}\|_\infty
          &\le64L^{-15/8+3\varepsilon/2}\theta_1,\\
 \|\partial\mathscr A^{\mathrm b}\|_\infty
          &\le137L^{-15/8+3\varepsilon/2}\theta_2,\\
 \|\delta_\alpha\partial\mathscr A^{\mathrm b}\|_\infty
          &\le1382L^{-15/8+3\varepsilon/2}\theta_3,\\
 \|\mathscr C^{\mathrm b}\|_\infty&<L^{-7/4}d_1,\qquad
 \|\partial\mathscr C^{\mathrm b}\|_\infty<L^{-7/4}d_2,
 \end{split}
\end{equation}
and
\begin{equation}\label{eq:rg-small-polymer-constant-bounds}
 \|\mathscr K_1\|_\infty<b_1,\qquad
 \|\mathscr K_2\|_\infty<L^{-1+\varepsilon}d_1.
\end{equation}
Consequently
\begin{equation}\label{eq:rg-small-polymer-analytic-disk}
 (t\mathscr A^{\mathrm b},t\mathscr C^{\mathrm b})
                       \in R_k\times\mathfrak X_k
                             \qquad(|t|\le L^{7/4}).
\end{equation}
The construction uses a base point of $X$ and a decomposition
and local real representative of $\mathscr A$, and is
independent of the activity. When $\mathscr C=0$, it gives
$\mathscr C^{\mathrm b}=\mathscr K_2=0$.
\end{proposition}

\begin{proof}
Choose a lift of $X$ and a photon site $c_*$ in it.
Set $r=2LM\ge50$ and $K_{\mathrm{loc}}=K(c_*,2r+5)$, with
sup balls understood through their nonwrapping lifts.
Choose the coarse $M$-cube containing $c_*/L$.
If $\mathscr A=A_0+A_1$ is its domain decomposition, its
real representative $A'$ and gauge $\lambda$ on the
enlarged cube give
\[
 a_0=(A_0)_L=a'_0+d\lambda_L,\qquad
 a'_0=(A')_L,\qquad a_1=(A_1)_L
\]
on $K_{\mathrm{loc}}$. To see that these functions are
defined there, the scaled enlarged cube has side
$(2L+1)LM$ and center $z$ with $|c_*-z|<LM/2$.
The inequality $4LM+6\le L^2M$ puts
$K(c_*,2r+6)$ inside this cube. Moreover
$4r+11<\ell$, since $\ell\ge13L^2M$.
Thus the boxes and all derivative cells used below do
not wrap.

We first record the cutoff estimate with its derivative
support. Put
\[
 f(t)=
 \begin{cases}
 1,&|t|\le1,\\
 1-2(|t|-1)^2,&1\le|t|\le3/2,\\
 2(2-|t|)^2,&3/2\le|t|\le2,\\
 0,&|t|\ge2.
 \end{cases}
\]
For a bond $(c,c+\eta_ke_\mu)$ use the product of the three
functions $f((c_\nu+\eta_k\delta_{\mu\nu}/2-c_{*,\nu})/r)$
as its cutoff multiplier. Denote this operation by
$\operatorname{cut}_r$. The function $f$ is continuously
differentiable, $|f'|\le2$, and $f'$ is Lipschitz with
constant $4$. If a bond function $b$ has local norms
$B_0,B_1,B_\alpha$ on $K_{\mathrm{loc}}$, extension by zero
after multiplication gives
\begin{equation}\label{eq:rg-cutoff-estimate}
 \begin{split}
 \|\operatorname{cut}_r b\|_\infty&\le B_0,\\
 \|\partial\operatorname{cut}_r b\|_\infty
                                      &\le B_1+2B_0/r,\\
 \|\delta_\alpha\partial\operatorname{cut}_r b\|_\infty
                             &\le B_\alpha+12B_1/r+12B_0/r^2.
 \end{split}
\end{equation}
Indeed, the multiplier has oscillation at most
$6|c-c'|/r$, first difference at most $2/r$, and
first-difference oscillation at most $12|c-c'|/r^2$.
Apply the product rule to the first difference.
In its H\"older difference the four terms are bounded
by $B_\alpha |c-c'|^\alpha$, twice
$6B_1|c-c'|/r$, and $12B_0|c-c'|/r^2$.
Here $|c-c'|\le1$ implies
$|c-c'|\le|c-c'|^\alpha$.
If any such term is nonzero, all its cells lie in
$K(c_*,2r+5)$: a nonzero multiplier or its difference
requires a site within $2r+1$, the first-difference
cell adds one mesh length, and the pair adds at most
$1$. The same bound covers the two-step cell at
$k=0$. Thus the local estimates suffice for the global
ones. The cutoff equals one on bonds with both ends in
$K(c_*,r)$.

For $b=a'_0$ or $a_1$, write $v=c-c_*$ and define
\[
 V[b]_{\nu\mu}=(\partial_\nu b_\mu)(c_*),\qquad
 \operatorname{Lin}[b]_\mu(c)=\sum_\nu V[b]_{\nu\mu}v_\nu,
 \qquad
 \operatorname{rem}[b]=b-b(c_*)-\operatorname{Lin}[b].
\]
Its first difference is
$\partial_\nu b_\mu(c)-\partial_\nu b_\mu(c_*)$,
and its H\"older first-difference seminorm is at most
$b_3$. Join the two sites by at most $2r+5$ simultaneous
coordinate steps of length at most $1$. Their difference
is therefore at most $(2r+5)b_3\le3rb_3$. A rectilinear
path to $c_*$ has length at most $3(2r+5)$, and the
remainder vanishes at $c_*$. Keeping the sharper derivative
bound along this path and using $r\ge50$ gives
\[
 \|\operatorname{rem}[b]\|_{K_{\mathrm{loc}}}
              \le3(2r+5)^2b_3\le18r^2b_3,\qquad
 \|\partial\operatorname{rem}[b]\|_{K_{\mathrm{loc}}}
              \le3rb_3.
\]
Equation \eqref{eq:rg-cutoff-estimate} consequently bounds
the three global norms of its cutoff by
$18r^2b_3$, $39rb_3$, and $253b_3$, respectively.

Only a real gauge is available. We remove the constant
and symmetric linear part of $a'_0$, leaving its curvature
and the complex linear part of $a_1$. Set
\[
 V^0=V[a'_0],\qquad
 V^{\mathrm s}=(V^0+(V^0)^T)/2,\qquad
 F=V^0-(V^0)^T,
\]
and use the real quadratic function
\begin{equation}\label{eq:rg-real-quadratic-gauge}
 q_*(c)=\sum_\mu\left[a'_{0,\mu}(c_*)-
                    \frac{\eta_k}{2}V^{\mathrm s}_{\mu\mu}\right]v_\mu
              +\frac12\sum_{\nu,\mu}V^{\mathrm s}_{\nu\mu}v_\nu v_\mu.
\end{equation}
Its finite difference is exactly
$(dq_*)_\mu(c)=a'_{0,\mu}(c_*)+
\sum_\nu V^{\mathrm s}_{\nu\mu}v_\nu$.
The mesh term in \eqref{eq:rg-real-quadratic-gauge} cancels
the diagonal term from differencing the quadratic polynomial.
With $\mathscr K_1=a_1(c_*)$ and
\[
 P_\mu(c)=\sum_\nu
                    \left[\frac12F_{\nu\mu}+V[a_1]_{\nu\mu}\right]v_\nu,
\]
we have the local identity
\begin{equation}\label{eq:rg-local-taylor-gauge}
 \mathscr A_L-d(\lambda_L+q_*)=
       \mathscr K_1+P+\operatorname{rem}[a'_0]+\operatorname{rem}[a_1].
\end{equation}

To cut off $P$ we need a bound that also uses the side
$\ell$. A global mean-zero periodic function whose
oscillation between sites at distance at most $1$ is at
most $B$ has sup norm at most $\ell B$. Any two sites
can be joined by at most $\lceil\ell/2\rceil\le\ell$
simultaneous coordinate steps of length at most $1$;
comparison with the mean gives the claim.
Every $\partial_\nu a_{1,\mu}$ has mean zero, sup norm
at most $b_2$, and this oscillation at most $b_3$.
For $a_0$, the global curl agrees on each derivative
cell with the curl of a local real representative.
Its sup norm is at most $2b_2$, its oscillation at most
$2b_3$, and it has mean zero. The cells of a pair at
distance at most $1$ lie in one enlarged cube, so that
the same representative can be used for their comparison.
It follows that every coefficient of $P$ is bounded by
\begin{equation}\label{eq:rg-linear-slope-bound}
 s=2\min(b_2,\ell b_3).
\end{equation}

Use a second radius
\[
 r'=\min(e_k^{-\varepsilon},\ell/5).
\]
The charge condition in \eqref{eq:rg-contraction-parameters}
gives $r'\ge r$, and $4r'+11<\ell$. Thus its cutoff is
also defined without wrapping. The three norms of $P$
on its support are at most $6.3r's$, $s$, and $0$.
Equation \eqref{eq:rg-cutoff-estimate}, with $r'$ in
place of $r$, bounds the cutoff norms by
$6.3r's$, $13.6s$, and $87.6s/r'$.
If $r'=e_k^{-\varepsilon}$, use $s\le2b_2$;
in units of $(\theta_1,\theta_2,\theta_3)$ these are
at most $(12.6,27.2,175.2)L^{-15/8+\varepsilon}$.
If $r'=\ell/5<e_k^{-\varepsilon}$, use
$s\le2\ell b_3$. Since $\ell e_k^\varepsilon<5$,
the corresponding three bounds are at most
$(63,136,876)L^{-15/8+3\varepsilon/2}$.

Define the global field and real gauge by
\begin{equation}\label{eq:rg-global-representative-construction}
 \begin{split}
 \mathscr A^{\mathrm b}
   &=\operatorname{cut}_r
             \big(\operatorname{rem}[a'_0]+\operatorname{rem}[a_1]\big)
                  +\operatorname{cut}_{r'}P,\\
 \lambda_*&=\lambda_L+q_*\quad\text{on }K_{\mathrm{loc}},
                  \qquad \lambda_*=0\quad\text{elsewhere}.
 \end{split}
\end{equation}
The two remainder cutoffs contribute
$(36r^2,78r,506)b_3$. As $r=2LM$, the charge condition
$160L^{1+3\varepsilon/2}Me_k^\varepsilon\le1$ bounds
the first two by $L^{-15/8+3\varepsilon/2}\theta_1$
and $L^{-15/8+3\varepsilon/2}\theta_2$, respectively.
Adding the linear contribution gives the constants
$64$, $137$, and $1382$ in
\eqref{eq:rg-small-polymer-global-bounds}.

In the external argument take
\[
 \mathscr K_2=\mathscr C_L(c_*),\qquad
 \mathscr C^{\mathrm b}
                   =\operatorname{cut}_r(\mathscr C_L-\mathscr K_2).
\]
A rectilinear path to $c_*$ bounds the undifferentiated
field on $K_{\mathrm{loc}}$ by
$6.3r\|\partial\mathscr C_L\|_\infty$.
The cutoff bounds therefore give
\[
 \|\mathscr C^{\mathrm b}\|_\infty
       \le6.3r\|\partial\mathscr C_L\|_\infty,\qquad
 \|\partial\mathscr C^{\mathrm b}\|_\infty
       \le13.6\|\partial\mathscr C_L\|_\infty.
\]
Use \eqref{eq:rg-scaled-external-thresholds} and
$d_2=e_k^\varepsilon d_1$. The first charge condition
involving $160$ gives the desired bound for the sup norm;
the large-$L$ condition gives
$13.6L^{-2+3\varepsilon/2}\le L^{-7/4}$ for the derivative.
This proves the external bounds and the constant bounds.

The diameter estimate in Appendix~\ref{app:polymer-geometry}
puts $X$ in $K(c_*,M(L+2))$. Every bond with an endpoint
within distance $1$ of $X$ has both ends in $K(c_*,r)$,
since $M(L+2)+2\le2LM$. All cutoffs in the local identity
are one there, and the gauge is $\lambda_L+q_*$ there.
This proves \eqref{eq:rg-small-polymer-gauge-identity}.
Finally, the large-$L$ condition implies
$1382L^{-15/8+3\varepsilon/2}<L^{-7/4}$.
Multiplication by any $|t|\le L^{7/4}$ consequently leaves
all five norms strictly below their domain thresholds,
proving \eqref{eq:rg-small-polymer-analytic-disk}.
\end{proof}

\begin{proposition}\label{prop:rg-small-polymer-kernel-comparison}
Let $F(X)$ be a level-$k$ Grassmann activity holomorphic on
$\widetilde R_k\times\mathfrak X_k$, with the kernel support
and coefficient locality specified in
\S\ref{subsec:rg-activity-spaces}, and covariant under the
independent real gauges of its two connection arguments.
For the fields of Proposition~\ref{prop:rg-small-polymer-representatives},
its kernels satisfy
\begin{equation}\label{eq:rg-small-polymer-kernel-comparison}
 |F_{nm}(X;\mathscr A_L,\mathscr C_L;\xi;\zeta)|
   \le\exp(2(n+m))
        |F_{nm}(X;\mathscr A^{\mathrm b},\mathscr C^{\mathrm b};\xi;\zeta)|.
\end{equation}
For the vacuum kernel the two values are equal.
\end{proposition}

\begin{proof}
For each kernel argument $a$ let $c(a)$ be the photon site
of its ordinary leg, or the first photon site of its pair
leg, and let $\omega(a)\in\{0,1\}$ distinguish unbarred and
barred legs. On kernels the gauge homomorphism has the
form
\begin{equation}\label{eq:rg-kernel-gauge-phase}
 (G_\lambda F)_{nm}(\xi;\zeta)
       =F_{nm}(\xi;\zeta)
          \prod_a\exp\{(-1)^{\omega(a)}ie_k\lambda(c(a))\}.
\end{equation}
This follows by applying the homomorphism to the kernel
expansion; the multiplying product is symmetric in each
set of arguments, and hence preserves their antisymmetry.
The gauge relation \eqref{eq:rg-small-polymer-gauge-identity}
first gives
\[
 G_{\lambda_*}F(X;\mathscr A^{\mathrm b}+\mathscr K_1,
                         \mathscr C^{\mathrm b}+\mathscr K_2)
                       =F(X;\mathscr A_L,\mathscr C_L).
\]
Both pairs lie in the domains: a constant changes only
the sup norm, and
$L^{-7/4}+L^{-7/8+\varepsilon/2}<1$ and
$L^{-7/4}+L^{-1+\varepsilon}<1$.
Every phase from this real gauge has modulus one.

We remove the constant fields by continuation of a real
gauge identity. Write $\mathscr K_j=K_j^{\mathrm R}+iK_j^{\mathrm I}$
with real constant fields. On the connected domain
\[
 \Omega=\{z\in\C^4:|\operatorname{Im}z_1|<2,\,
       |\operatorname{Im}z_2|<2,\,|z_3|<2,\,|z_4|<2\},
\]
set
\[
 A(z)=\mathscr A^{\mathrm b}+z_1K_1^{\mathrm R}+z_2K_1^{\mathrm I},
 \qquad
 C(z)=\mathscr C^{\mathrm b}+z_3K_2^{\mathrm R}+z_4K_2^{\mathrm I},
\]
and set
\[
 \Lambda_z(c)=-
       (z_1K_1^{\mathrm R}+z_2K_1^{\mathrm I}
                  +z_3K_2^{\mathrm R}+z_4K_2^{\mathrm I})\cdot(c-c_*)
                    \quad(c\in K(c_*,r)),
\]
with value zero elsewhere. The first connection belongs to
$\widetilde R_k$: its complex small summand has sup norm
below $(L^{-7/4}+4L^{-7/8+\varepsilon/2})\theta_1<\theta_1$
and derivatives those of $\mathscr A^{\mathrm b}$, and its
remaining constant summand is real. Invariance under real
constant additions handles its unrestricted real part.
The external connection has sup norm below
$(L^{-7/4}+4L^{-1+\varepsilon})d_1<d_1$, and derivative
that of $\mathscr C^{\mathrm b}$. The large-$L$ condition
implies these two strict inequalities.

For fixed kernel arguments the function
\[
 H(z)=F_{nm}(X;A(z),C(z);\xi;\zeta)
             \prod_a\exp\{(-1)^{\omega(a)}ie_k\Lambda_z(c(a))\}
\]
is holomorphic on $\Omega$. When $z$ is real, split
$\Lambda_z$ into its first two and last two terms.
These are independent real gauge functions for the two
connections, with differentials equal to the negatives
of their respective constants on $K(c_*,r)$. Every bond
read by the local gauge condition lies in this ball.
The real gauge identity and
\eqref{eq:rg-kernel-gauge-phase} therefore give
$H(z)=H(0)$ for $z\in\R^2\times(-2,2)^2$.
All derivatives of $H-H(0)$ at zero vanish, so its
holomorphic power series is zero there. The identity
theorem extends this equality to the connected domain
$\Omega$. Evaluating at $z=(1,i,1,i)$ gives the identity
for the two original complex constants.

It remains to bound the inverse phases. By locality,
an ordinary argument lies over $X$. For a pair of the
first kind its first photon site also lies in $X$.
For a boundary pair that site has distance at most
$\eta_k\le1$ from $X$. Hence in all nonzero kernel terms
\[
 |c(a)-c_*|_1\le3(M(L+2)+1)\le9LM.
\]
The constant bounds and the charge conditions give
\[
 \begin{split}
 e_k|K_1^{\mathrm I}\cdot(c(a)-c_*)|
       &\le9LM L^{-7/8+\varepsilon/2}e_k^{1/4+\varepsilon}
                                                    \le1,\\
 e_k|K_2^{\mathrm I}\cdot(c(a)-c_*)|
       &\le9ML^\varepsilon e_k^{2\varepsilon}\le1.
 \end{split}
\]
Each leg phase, and its inverse, thus has modulus at most
$\exp(2)$. Combining this bound with the unit modulus
phases from $\lambda_*$ proves
\eqref{eq:rg-small-polymer-kernel-comparison}.
The vacuum has no leg phases, which proves its equality.
\end{proof}
\section{Contraction of the normalized activity}
\label{subsec:rg-activity-contraction}

On a small polymer the ordinary quadratic kernel now vanishes
at zero connection. The vacuum is also even in the connection,
so its first derivative vanishes there. The analytic
disk gives an additional gain for precisely these two
terms; every other term is controlled directly by its
Grassmann degree.

\begin{proposition}\label{prop:rg-activity-contraction}
Assume \eqref{eq:rg-adjustment-charge-condition},
\eqref{eq:rg-contraction-parameters}, and \eqref{eq:rg-contraction-polymer-parameters}.
For $\mathscr E\in K_k^{x,\C}(\ell)$ set
$E=(R\mathscr E)^{\mathrm{nat}}$.
The linear map $\mathscr E\mapsto\mathcal L E$ takes
$K_k^{x,\C}(\ell)$ into $K_{k+1}^{x,\C}(\ell/L)$,
preserves the real subspaces, and satisfies
\begin{equation}\label{eq:rg-contraction-norm}
 \|\mathcal L E\|_{k+1}^x
       \le c_{\mathrm{fl}}L^{-1/4+3\varepsilon/2}\|\mathscr E\|_k^x
       \le c_{\mathrm{fl}}L^{-1/4+2\varepsilon}\|\mathscr E\|_k^x,
       \qquad c_{\mathrm{fl}}=61c_{\mathrm{rb}}.
\end{equation}
For the weights $w_j^{\mathrm E}=e_j^{1/4-7\varepsilon}$,
\begin{equation}\label{eq:rg-contraction-window}
 \frac{\|\mathcal L E\|_{k+1}^x}{w_{k+1}^{\mathrm E}}
       \le c_{\mathrm{fl}}L^{-3/8+11\varepsilon/2}
                       \frac{\|\mathscr E\|_k^x}{w_k^{\mathrm E}}.
\end{equation}
Restriction to the zero external argument commutes with
the map.
\end{proposition}

\begin{proof}
Write
\[
 N(X)=\sup_{\widetilde R_k\times\mathfrak X_k}
                                      \|E(X;\cdot,\cdot)\|_{h_k}.
\]
Extraction and adjustment give
\begin{equation}\label{eq:rg-adjusted-polymer-bound}
 N(X)\le
 \begin{cases}
 20\|\mathscr E\|_k^x\exp(-\kappa d_M(X)),&X\text{ small},\\
 \|\mathscr E\|_k^x\exp(-\kappa d_M(X)),&X\text{ large}.
 \end{cases}
\end{equation}
For a small polymer this uses
$\|R\mathscr E(X;\mathscr A,\mathscr C)\|_{h_k}
\le\|\mathscr E(X;\mathscr A,\mathscr C)\|_{h_k}
 +\|\mathscr E(X;0,0)\|_{h_k}$ and
\eqref{eq:rg-adjustment-norm}. For a large polymer both
operations leave the activity unchanged.

Fix a small $X$ and a coarse connection pair, and use
the global representatives above. Put $\rho=L^{7/4}$.
The inverse image of the open set $R_k\times\mathfrak X_k$
under $t\mapsto(t\mathscr A^{\mathrm b},t\mathscr C^{\mathrm b})$
contains the closed disk $|t|\le\rho$, hence a slightly
larger disk. The functions
\[
 f_{nm}(t)=E_{nm}(X;t\mathscr A^{\mathrm b},
                              t\mathscr C^{\mathrm b})
\]
are holomorphic there, with values in a finite-dimensional
kernel space. Their weighted norms obey
\[
 \frac{h_{k,1}^nh_{k,2}^m}{n!m!}\|f_{nm}(t)\|_1
                                      \le N(X)\qquad(|t|\le\rho).
\]
Here $\|\cdot\|_1$ is the sum of kernel moduli with their
index weights.

For the vacuum, \eqref{eq:rg-adjusted-kernel-normalization}
gives $f_{00}(0)=0$. Charge conjugation fixes a scalar
kernel and reverses both connections, so
$f_{00}(-t)=f_{00}(t)$ and $f_{00}'(0)=0$.
Cauchy's formula gives
\[
 f_{00}(1)=\frac1{2\pi i}\oint_{|t|=\rho}
                                   \frac{f_{00}(t)}{t^2(t-1)}\,\dd t.
\]
The vacuum equality in
Proposition~\ref{prop:rg-small-polymer-kernel-comparison}
and $\rho\ge2$ therefore give
\begin{equation}\label{eq:rg-vacuum-cauchy-gain}
 \|(S_LE)^{[0,0]}(X;\mathscr A_L,\mathscr C_L)\|_{h_{k+1}}
              \le\frac{N(X)}{\rho(\rho-1)}
              \le2L^{-7/2}N(X).
\end{equation}

For the ordinary quadratic kernel, normalization gives
$f_{20}(0)=0$ as a full antisymmetric kernel. Its balanced
matrix components vanish by
\eqref{eq:rg-adjusted-kernel-normalization}; their
oppositely ordered components vanish by antisymmetry,
and the two equally barred or equally unbarred components
vanish by balance. Hence
\[
 f_{20}(1)=\frac1{2\pi i}\oint_{|t|=\rho}
                                    \frac{f_{20}(t)}{t(t-1)}\,\dd t,
 \qquad
 \|f_{20}(1)\|_1\le\frac{\sup_{|t|=\rho}\|f_{20}(t)\|_1}{\rho-1}.
\]
The comparison of kernels costs $\exp(4)$ in this
bidegree. Its scaling exponent is $p(2,0)=9/4$.
Multiplying the last inequality by $h_{k,1}^2/2$ and
using \eqref{eq:rg-degree-scaling} gives
\begin{equation}\label{eq:rg-quadratic-cauchy-gain}
 \|(S_LE)^{[2,0]}(X;\mathscr A_L,\mathscr C_L)\|_{h_{k+1}}
                       \le2\exp(4)L^{-4}N(X).
\end{equation}

For all remaining bidegrees use the original scaled
connection pair, which belongs to the fine domain.
Evenness leaves only $(1,1)$, $(0,2)$, and total degrees
at least $4$. The exponents satisfy
\[
 p(1,1)>\frac{13}{4}-\frac{3\varepsilon}{2},\qquad
 p(0,2)>\frac{13}{4}-\frac{3\varepsilon}{2},\qquad
 p(n,m)\ge n+m.
\]
Consequently
\[
 \begin{split}
 \sum_{(n,m)\ne(0,0),(2,0)}
       \|(S_LE)^{[n,m]}(X;\mathscr A_L,\mathscr C_L)\|_{h_{k+1}}
 &\le L^{-13/4+3\varepsilon/2}
       \sum_{n,m}\|E^{[n,m]}(X;\mathscr A_L,\mathscr C_L)\|_{h_k}\\
 &\le L^{-13/4+3\varepsilon/2}N(X).
 \end{split}
\]
The sum is taken at one connection pair before bounding
it by $N(X)$. The large-$L$ condition ensures
$2L^{-7/2}\le L^{-13/4}$ and
$2\exp(4)L^{-4}\le L^{-13/4}$. Adding the three
degree contributions yields
\begin{equation}\label{eq:rg-small-polymer-contraction}
 \|S_LE(X;\mathscr A_L,\mathscr C_L)\|_{h_{k+1}}
                       \le3L^{-13/4+3\varepsilon/2}N(X).
\end{equation}

We now sum over the fine polymers in a coarse polymer
$Y$. The hypotheses \eqref{eq:rg-contraction-polymer-parameters}
imply
$\kappa-1\ge\kappa_0$ and
$L(\kappa-\kappa_0-2)\ge\kappa$.
The reblocking bound
\eqref{eq:polymer-reblock-sum} at rate $\kappa$
and \eqref{eq:rg-adjusted-polymer-bound} therefore give
\[
 \begin{split}
 \sum_{\substack{X\text{ small}\\\bar X^L=LY}}
       \|S_LE(X;\mathscr A_L,\mathscr C_L)\|_{h_{k+1}}
 &\le60c_{\mathrm{rb}}L^{-1/4+3\varepsilon/2}
       \|\mathscr E\|_k^x\exp(-\kappa d_M(Y)).
 \end{split}
\]
For large $X$, scaling decreases the norm by
\eqref{eq:rg-scaling-norm-identity} and the ratios of
the two weights. Since $d_M(X)>L$,
\[
 \exp(-\kappa d_M(X))
       \le\exp(-L)\exp(-(\kappa-1)d_M(X)).
\]
Use \eqref{eq:polymer-reblock-sum} at rate $\kappa-1$
to obtain
\[
 \sum_{\substack{X\text{ large}\\\bar X^L=LY}}
       \|S_LE(X;\mathscr A_L,\mathscr C_L)\|_{h_{k+1}}
 \le c_{\mathrm{rb}}L^3\exp(-L)
             \|\mathscr E\|_k^x\exp(-\kappa d_M(Y)).
\]
As $L^3\exp(-L)\le L^{-1/4}$, these estimates prove
\eqref{eq:rg-contraction-norm} after taking the two
suprema in the activity norm. Division by
$w_{k+1}^{\mathrm E}=L^{1/8-7\varepsilon/2}w_k^{\mathrm E}$
gives the relative factor $c_{\mathrm{fl}}L^{-3/8+5\varepsilon}$,
which is at most the factor in \eqref{eq:rg-contraction-window}.
The additional parameter choice
$c_{\mathrm{fl}}L^{-3/8+6\bar\varepsilon}\le1/16$
makes its relative factor at most $1/16$.

The map is linear. Its kernels are finite sums of
holomorphic fine kernels composed with the complex
linear connection scaling. The locality was proved
after \eqref{eq:rg-reblocking-definition}.
For lattice covariance a coarse group element
$g'=(S,a')$ has the fine representative $g=(S,La')$.
It gives a bijection of the polymers in the two
reblocking sums, satisfies
$(g'\mathscr A)_L=g\mathscr A_L$, and intertwines the
spin actions by $a_{g'}S_L=S_La_g$.
The same identity holds for $\mathscr C$.
Thus covariance of each fine activity gives covariance
of its sum. For the real gauges, every fine bond within
distance $1$ of $X$ scales to a bond within distance
$1/L\le1$ of $Y$, while
$G_\lambda S_L=S_LG_{\lambda_L}$ follows from
$e_k\lambda_L(c)=e_{k+1}\lambda(c/L)$.
Independent real gauges of the two arguments therefore
pass through each term. Finally, scaling has real
coefficients, preserves the bar index, and commutes with
charge conjugation and the antilinear reality map.
These facts prove membership in both stated activity
spaces. Since extraction and adjustment commute with
restriction and $0_L=0$, reblocking also commutes with
restriction to the zero external argument.
\end{proof}

The reblocked sum retains the evaluated normalized
activity. Indeed, each fine polymer belongs to exactly
one reblocking class, so
\begin{equation}\label{eq:rg-reblocked-sum}
 \sum_Y(\mathcal LE)(Y;\mathscr A,\mathscr C)
      =S_L\sum_X E(X;\mathscr A_L,\mathscr C_L).
\end{equation}
If the coarse generators are evaluated at fields
with $\chi'=\delta_{\alpha,\mathscr A+\mathscr C}\psi'$
and its barred counterpart, their scaled fields satisfy
the fine relations. This follows from
$d_X(x/L,x'/L)=d_X(x,x')/L$ and
$U_{(\mathscr A+\mathscr C)_L}(x,x')
 =U_{\mathscr A+\mathscr C}(x/L,x'/L)$:
the path bijection and
$e_k\eta_kL^{-1/2}=e_{k+1}\eta_{k+1}$ give the latter.
Substitution in the difference quotient produces
$L^{-1-\alpha}\chi'(x/L,x'/L)$, and the barred quotient
uses the inverse transport. Adjustment consequently
has its original value under this evaluation. Thus the
evaluation of \eqref{eq:rg-reblocked-sum} is also the
evaluation of
\[
 S_L\sum_X(R\mathscr E)(X;\mathscr A_L,\mathscr C_L).
\]
\begin{proof}[Proof of Theorem~\ref{thm:main}]
Proposition~\ref{prop:rg-scalar-counterterms} gives the scalar
counterterms and their bounds. Equations
\eqref{eq:rg-exact-extraction-identity} and
\eqref{eq:rg-remainder-properties} give the exact extraction and
the integrated normalization. The shortest-chain construction
\eqref{eq:rg-chain-homomorphism} preserves evaluated values;
\eqref{eq:rg-adjusted-kernel-normalization} and
\eqref{eq:rg-adjustment-norm} give the full kernel normalization
and its bound. Proposition~\ref{prop:rg-activity-contraction}
proves the contraction and the relative estimate. The covariance
and restriction statements were proved with each map.
\end{proof}
\section{Local reconstruction of a mean field}
\label{subsec:rg-field-reconstruction}

The result in this section follows
\cite[\S\S12.4,13.4]{QED3Manuscript} and has hypotheses independent
of the large-\(L\) contraction conditions. Fix \(k\ge1\),
\(\eta_k=L^{-k}\), \(L\ge5\) odd, \(M=L^m\) with \(m\in\Z_{\ge2}\),
and an admissible torus as in Section~\ref{sec:activities}.
Let \(0<\alpha<1\), \(0<\rho<1\), \(0<\mu\le1\), and
\(b_k\ge b_->0\).

The chamber Dirac stencil is
\[
 D_A\psi(x)=-\eta_k^{-1}\sum_{i=0}^3M_i(x)
            [U_A(x,x^{(i)})\psi(x^{(i)})-\psi(x)],
 \qquad
 M_i(x)=\tfrac12w_i(\rho-\gamma\cdot\widehat d_i(x)),
\]
where
\[
 \begin{gathered}
 w_0=w_3=6,\qquad w_1=w_2=6\sqrt2,\\
 \widehat d_0=-ge_3,\quad
 \widehat d_1=-g(e_2-e_3)/\sqrt2,\quad
 \widehat d_2=-g(e_1-e_2)/\sqrt2,\quad
 \widehat d_3=ge_1 .
 \end{gathered}
\]
In particular
\(\sum_iw_i\widehat d_i=0\) and \(\|M_3(x)\|_2\le6\).

Let \(\Delta_y\) be the unit block centered at
\(y\in\Z^3/\ell\Z^3\). For \(c\in\Delta_y\), join \(c\) to the
successive centers of its nested blocks of sides
\(\eta_k,L\eta_k,\ldots,1\). Let \(\varphi_A(y,c)\) be the
exponential of \(ie_k\eta_k\) times the sum of the six-path
averages along these successive center-to-site paths, as in
\eqref{eq:rg-reconstruction-phase}. Define the normalized average
and its reverse by
\begin{equation}\label{eq:reconstruction-average}
 (\mathcal Q_k(A)v)(y)
      =\frac{\eta_k^3}{48}\sum_{c\in\Delta_y,g}
                                     \varphi_A(y,c)v(c,g),\qquad
 (\mathcal Q_k^\sharp(A)u)(c,g)=\varphi_A(y,c)^{-1}u(y).
\end{equation}
Thus \(\mathcal Q_k\mathcal Q_k^\sharp=1\).
Equivalently \(\mathcal Q_k^\sharp(A)=\mathcal Q_k(-A)^T\),
where \(T\) is the algebraic transpose with the stated fine
measure and unit counting measure. Reversing the connection and
taking this weighted transpose are separate operations.
Put
\[
 F_\infty(A)=D_A+\mu+b_k\mathcal Q_k^\sharp(A)\mathcal Q_k(A),
 \qquad
 T_k(A)=\mathcal Q_k(A)(b_k^{-1}D_A+1+\mu/b_k).
\]
Write \(\bar T_k(A)=\mathsf C^{-1}T_k(-A)\mathsf C\).
For a complex connection set
\(\mathfrak i(A)=e_k\|\operatorname{Im}A\|_\infty\) and
\(q_d(A)=e_k\|dA\|_\infty\).
For a scalar-component kernel \(P(a,b)\) acting with column
measure \(\nu(b)\), write
\[
 |P|_{1,\infty}=\sup_a\sum_b\nu(b)|P(a,b)|.
\]
The sum includes spin indices. For ordinary and pair columns
the measures are \(\eta_k^3/48\) and \((\eta_k^3/48)^2\),
respectively.

\begin{theorem}[Local reconstruction]\label{thm:reconstruction}
Under the preceding hypotheses the following statements hold.
\begin{enumerate}[label=\textup{(\roman*)},leftmargin=*]
\item For every complex connection \(A\) and every union \(Y\) of
partition \(M\)-cubes, the rows \(T_Y^{(1)}(A)\) and
\(T_Y^{(2)}(A)\) in \eqref{eq:rg-reconstruction-two-pieces} exist.
The first uses only ordinary legs over \(Y\); the second uses
the external ordered pairs leaving \(Y\). Their coefficients
read only bonds meeting \(Y\), are entire in \(A\), and satisfy
\[
 T_Y^{(1)}(A)\psi+T_Y^{(2)}(A)\delta_{\alpha,A}\psi
                       =T_k(A)\psi
 \quad\text{at the unit sites in }Y.
\]
They obey the lattice, real-gauge and reality covariance
identities \eqref{eq:rg-reconstruction-gauge-covariance}--%
\eqref{eq:rg-reconstruction-reality}. Their barred counterparts
are \(\bar T_Y^{(i)}(A)=\mathsf C^{-1}T_Y^{(i)}(-A)\mathsf C\).

\item In Euclidean spinor supremum norm, writing
\(\mathfrak i=\mathfrak i(A)\), \(q_d=q_d(A)\),
\begin{equation}\label{eq:rg-reconstruction-sup-estimate}
 \begin{aligned}
 \|T_Y^{(1)}(A)\|_{\mathrm{su}},\ \|T_k(A)\|_{\mathrm{su}}
 &\le e^{3\mathfrak i/2}
       [1+b_-^{-1}(54Lq_de^{6q_d}+12+6e^{\mathfrak i}+\mu)],\\
 \|T_Y^{(2)}(A)\|_{\mathrm{su}}
 &\le6b_-^{-1}e^{3\mathfrak i/2}\eta_k^\alpha .
 \end{aligned}
\end{equation}
If \(\mathfrak i,q_d\le1\), the coefficient row sums satisfy
\begin{equation}\label{eq:rg-reconstruction-coefficient-bound}
 |T_Y^{(i)}(A)|_{1,\infty},\ |T_k(A)|_{1,\infty},\
 |\bar T_Y^{(i)}(A)|_{1,\infty},\ |\bar T_k(A)|_{1,\infty}
 \le C_T\quad(i=1,2),
\end{equation}
where
\[
 C_T=2e^{3/2}[1+b_-^{-1}(54Le^6+12e+1)].
\]
This constant is independent of \(k,\ell,Y\). In particular
the bound holds for \(A\in\widetilde R_k+\mathfrak X_k\)
when \(0<\varepsilon\le1/200\), \(0<e_k<1\), and
\eqref{eq:rg-domain-charge-conditions} holds.

\item Introduce independent unit generators \(\Psi,\bar\Psi\)
over \(Y\). Replacing them by
\(T_Y^{(1)}\psi+T_Y^{(2)}\chi\) and its barred counterpart,
and fixing the old generators, defines the homomorphism
\(\operatorname{sub}_{Y,A}\) of
\eqref{eq:rg-step-local-substitution}. It preserves balance
and the symmetry identities
\eqref{eq:rg-step-substitution-covariances}.
If the original ordinary legs lie over \(Y\) and its pair legs
are admissible for \(Y\), the substituted polynomial has the same
support properties.
If \(\mathfrak i,q_d\le1\) and the ordinary, pair and unit
weights \(w_1,w_2,w_3>0\) satisfy
\(C_T(w_1+w_2)\le w_3\), then
\[
 \|\operatorname{sub}_{Y,A}F\|_{w_1,w_2}
                       \le\|F\|_{w_1,w_2;w_3}
\]
for every polynomial whose unit generators lie over \(Y\).

\item At each \(A\) for which \(F_\infty(A)\) is invertible,
define
\[
 H_k(A)=b_kF_\infty(A)^{-1}\mathcal Q_k^\sharp(A).
\]
Then \(T_k(A)H_k(A)=1\).
For the barred recovery also assume \(F_\infty(-A)\) is
invertible and define
\(\bar H_k(A)=\mathsf C^{-1}H_k(-A)\mathsf C\).
On evaluation at \(\psi=H_k(A)\Phi\),
\(\chi=\delta_{\alpha,A}\psi\), and the barred analogues,
the substitution sends \(\Psi,\bar\Psi\) to \(\Phi,\bar\Phi\)
over \(Y\).
\end{enumerate}
\end{theorem}

The inverse hypothesis in the last clause is additional to the
row estimate. The proof below constructs the rows directly.
Reindexing the internal chamber hops cancels the bulk part of
the Dirac stencil. The remaining boundary rows are controlled by
face counts, while the averaging-phase defect is controlled by
curvature and the first common block scale.

\subsection{The rows and their cancellation}
The normalized average \eqref{eq:reconstruction-average} gives
\begin{equation}\label{eq:rg-reconstruction-definition}
 T_k(A)=b_k^{-1}\mathcal Q_k(A)
            \big(D_A+\mu+b_k\mathcal Q_k^\sharp(A)\mathcal Q_k(A)\big)
       =\mathcal Q_k(A)\big(b_k^{-1}D_A+1+\mu/b_k\big).
\end{equation}
The second form follows from
$\mathcal Q_k\mathcal Q_k^\sharp=1$.
At a connection where $F_\infty(A)$ is invertible, its mean map
$H_k(A)=b_kF_\infty(A)^{-1}\mathcal Q_k^\sharp(A)$ satisfies
\begin{equation}\label{eq:rg-reconstruction-left-inverse}
 T_k(A)H_k(A)
 =b_k^{-1}\mathcal Q_k(A)F_\infty(A)
               b_kF_\infty(A)^{-1}\mathcal Q_k^\sharp(A)=1.
\end{equation}
The barred identity follows where the inverse at $-A$ exists,
by conjugation with $\mathsf C$.

For $c\in\Delta_y$, let $c_j(c)$ be the center of its block
of side $L^j\eta_k$, $0\le j\le k$, so that $c_0(c)=c$ and
$c_k(c)=y$. The averaging phase is
\begin{equation}\label{eq:rg-reconstruction-phase}
 \varphi_A(y,c)=\exp\left(ie_k\eta_k
                 \sum_{j=0}^{k-1}(\tau A)(c_{j+1}(c),c_j(c))\right).
\end{equation}
The sum uses the average of the six rectilinear path sums.
We write $(\tau_kA)(y,c)$ for this sum.
All of these paths lie in $\Delta_y$ and their total length
in bonds is at most $\frac32(L^k-1)$. Thus
\begin{equation}\label{eq:rg-reconstruction-phase-moduli}
 \begin{aligned}
 |\varphi_A(y,c)^{\pm1}|&\le\exp(3\mathfrak i(A)/2),\\
 |u_A(c,c')^{\pm1}|&\le\exp(\eta_k\mathfrak i(A))
                                             \le\exp(\mathfrak i(A)).
 \end{aligned}
\end{equation}
for a photon bond, with $u_A(c,c')=\exp(ie_k\eta_k A(c,c'))$.

Write \(c_3=\pi x^{(3)}=c+\eta_kge_1\). The row form of
\eqref{eq:rg-reconstruction-definition} is
\begin{equation}\label{eq:rg-reconstruction-row}
 \begin{aligned}
 (T_k(A)\psi)(y)
 &=\int_{\Delta_y}\dd x\,\varphi_A(y,c)
                  \Bigl\{(1+\mu/b_k)\psi(x)\\
 &\qquad-(b_k\eta_k)^{-1}\sum_{i=0}^3M_i(x)
           [U_A(x,x^{(i)})\psi(x^{(i)})-\psi(x)]\Bigr\}.
 \end{aligned}
\end{equation}
Here $U_A=1$ for $i\le2$ and $U_A(x,x^{(3)})=u_A(c,c_3)$.
The normalized chamber integral over $\Delta_y$ has total mass
one. This row exists for every complex connection, regardless
of invertibility; the inverse hypothesis was needed only for
\eqref{eq:rg-reconstruction-left-inverse}.

The rows map antiperiodic fields to antiperiodic coarse fields:
translation by a period preserves blocks, neighbors and phases,
and supplies the single antiperiodic sign from the source field.

If $Y$ is a union of $M$-cubes and $y\in Y$ is a unit site,
then $\Delta_y\subset Y$: both $M$ and $L^k$ are odd, and the
unit block is contained in the $M$-cube of its center.
Define the two rows
\begin{equation}\label{eq:rg-reconstruction-two-pieces}
 \begin{split}
 (T_Y^{(1)}(A)\psi)(y)
 &=\int_{\Delta_y}\dd x\,\varphi_A(y,c)
    \Bigl\{(1+\mu/b_k)\psi(x)\\
 &\qquad-(b_k\eta_k)^{-1}\!\!\sum_{i:\,x^{(i)}\text{ over }Y}
           M_i(x)[U_A(x,x^{(i)})\psi(x^{(i)})-\psi(x)]\Bigr\},\\
 (T_Y^{(2)}(A)\chi)(y)
 &=-b_k^{-1}\eta_k^{\alpha-1}\int_{\Delta_y}\dd x\,
                1[c_3\notin Y]\varphi_A(y,c)M_3(x)\chi(x,x^{(3)}).
 \end{split}
\end{equation}
Rows outside $Y$ are zero. The first row contains ordinary
legs only over $Y$. The second row contains precisely the
type-$3$ pairs leaving $Y$, which are admissible boundary
pairs of length $\eta_k$. Their coefficients read $A$ only
on bonds with an endpoint in $\Delta_y$.
If $\chi=\delta_{\alpha,A}\psi$ on these boundary pairs, then
\begin{equation}\label{eq:rg-reconstruction-local-identity}
 T_Y^{(1)}(A)\psi+T_Y^{(2)}(A)\chi=T_k(A)\psi
                                             \quad\text{on the unit sites of }Y.
\end{equation}
This follows by replacing the missing type-$3$ difference in
\eqref{eq:rg-reconstruction-row} by
$\eta_k^\alpha\chi(x,x^{(3)})$.
The intra-cluster differences remain in $T_Y^{(1)}$.

\begin{proposition}[Uniform row bounds]\label{prop:rg-reconstruction-estimate}
The rows \eqref{eq:rg-reconstruction-two-pieces} and $T_k(A)$
obey the Euclidean norm estimates
\eqref{eq:rg-reconstruction-sup-estimate} for every complex
connection $A$. If $\mathfrak i(A),q_d(A)\le1$, their coefficient
row sums obey \eqref{eq:rg-reconstruction-coefficient-bound},
with the stated constant $C_T$ independent of $k,\ell,Y$.
\end{proposition}

\begin{proof}
Write $\mathfrak i=\mathfrak i(A)$ and $q_d=q_d(A)$.
Put $E=\Delta_y\times O_{\mathrm{cub}}$ and divide it into
\[
 E_1=\{x\in E:c_3\in\Delta_y\},\quad
 E_2=\{x\in E:c_3\in Y\setminus\Delta_y\},\quad
 E_3=\{x\in E:c_3\notin Y\}.
\]
For each coordinate direction $a\in\{\pm e_\mu\}$ there are
eight labels with $ge_1=a$. The face of $\Delta_y$ through
which $c+\eta_ka$ leaves the unit block has $\eta_k^{-2}$
photon sites. Thus
\begin{equation}\label{eq:rg-reconstruction-face-count}
 \#E=48\eta_k^{-3},\qquad \#(E_2\cup E_3)=48\eta_k^{-2}.
\end{equation}

The chamber hopping matrices satisfy
\[
 M_i(x^{(i)})-M_i(x)=w_i\gamma\cdot\widehat d_i(x),\qquad
 \sum_{i=0}^3w_i\gamma\cdot\widehat d_i(x)=0.
\]
Also $\|M_3(x)\|_2\le6$ and
$\|w_3\gamma\cdot\widehat d_3(x)\|_2=6$.
Reindex each intra-cluster sum by $x\mapsto x^{(i)}$ and
each type-$3$ sum in $E_1$ by $x\mapsto x^{(3)}$.
Both are involutions of their summation sets. The intra-cluster
sums become
$\sum_{x\in E}\varphi_A(y,c)\sum_{i\le2}
 (M_i(x^{(i)})-M_i(x))\psi(x)$, which equals
$-\sum_{x\in E}\varphi_A(y,c)w_3\gamma\cdot\widehat d_3(x)\psi(x)$.
The part over $E_1$ cancels with the corresponding matrix
difference from the reindexed type-$3$ sum. Thus
$(T_Y^{(1)}(A)\psi)(y)=R_0+R_1+R_2+R_3$, with
\begin{equation}\label{eq:rg-reconstruction-summation-parts}
 \begin{split}
 R_0&=(1+\mu/b_k)\frac{\eta_k^3}{48}
                         \sum_{x\in E}\varphi_A(y,c)\psi(x),\\
 R_1&=b_k^{-1}\frac{\eta_k^2}{48}
       \sum_{x\in E_2\cup E_3}
                     \varphi_A(y,c)w_3\gamma\cdot\widehat d_3(x)\psi(x),\\
 R_2&=-b_k^{-1}\frac{\eta_k^2}{48}\sum_{x\in E_2}
       \varphi_A(y,c)M_3(x)[u_A(c,c_3)\psi(x^{(3)})-\psi(x)],\\
 R_3&=-b_k^{-1}\frac{\eta_k^2}{48}\sum_{x\in E_1}
       [\varphi_A(y,c_3)u_A(c_3,c)-\varphi_A(y,c)]
                                             M_3(x^{(3)})\psi(x).
 \end{split}
\end{equation}
The face count and phase moduli bound the first three terms by
\[
 \exp(3\mathfrak i/2)\|\psi\|_{\mathrm{su}}
       \big[1+\mu/b_k+6b_k^{-1}
                           +6b_k^{-1}(1+\exp(\mathfrak i))\big].
\]
The fourth term requires the curvature estimate.

For neighboring photon sites $c,c'=c+\eta_ka$ in $\Delta_y$,
let $j(c,a)$ be the first scale at which their block centers
agree. The centers at all larger scales then agree as well.
For $j\ge2$, such a $c$ lies on a face of a block of side
$L^{j-1}\eta_k$. There are $L^{3(k-j+1)}$ such blocks inside
$\Delta_y$, with $L^{2(j-1)}$ sites on the indicated face.
For $j=1$ use all sites. We obtain
\begin{equation}\label{eq:rg-first-common-block-count}
 \#\{c:c,c+\eta_ka\in\Delta_y,\ j(c,a)=j\}
                         \le\eta_k^{-3}L^{-(j-1)}.
\end{equation}

Write
\[
 \begin{aligned}
 \varphi_A(y,c')u_A(c',c)-\varphi_A(y,c)
    &=\varphi_A(y,c)(\exp(v(c,a))-1),\\
 v(c,a)&=ie_k\eta_k
      [(\tau_kA)(y,c')+A(c',c)-(\tau_kA)(y,c)].
 \end{aligned}
\]
The contributions above scale $j(c,a)$ cancel in this bracket.
The remaining expression is the mean of the line sums over
closed paths from their common scale-$j$ center down to $c'$,
then across the bond to $c$, and back up the reversed paths.
All these loops lie in the scale-$j$ block, of side at most
$L^j\eta_k$, and have at most $3L^j$ bonds.
For a closed path of $n$ bonds in a box of side $s$, the axial
path estimate gives $|A(\Gamma)|\le2ns\|dA\|_\infty$.
To recall the reason, subtract the gauge obtained by line sums
from one corner in coordinate order $1,2,3$. The difference on
each bond is a sum over at most two strips of width one mesh
and length at most $s$. Its modulus is at most
$2s\|dA\|_\infty$, and the gauge sum around the loop vanishes.
This argument applies to complex fields as well.
Consequently, with $j=j(c,a)$,
\begin{equation}\label{eq:rg-reconstruction-phase-defect}
 |v(c,a)|\le6L^{2(j-k)}q_d,\qquad
 |\exp(v(c,a))-1|
                    \le6L^{2(j-k)}q_d\exp(6q_d).
\end{equation}
The factor $L^{2(j-k)}$ must be kept until after summing the
locations of the first common block. By
\eqref{eq:rg-first-common-block-count},
\begin{align*}
 \eta_k^2\sum_{\substack{c\in\Delta_y\\c+\eta_ka\in\Delta_y}}
                                  |\exp(v(c,a))-1|
 &\le6q_d\exp(6q_d)L^{1-k}\sum_{j=1}^kL^j\\
 &\le9Lq_d\exp(6q_d).
\end{align*}
The six directions, eight labels per direction, and
$\|M_3\|_2\le6$ now give
\[
 |R_3|_2\le54Lq_d\exp(6q_d)b_k^{-1}
                         \exp(3\mathfrak i/2)\|\psi\|_{\mathrm{su}}.
\]
Together with the first three terms this proves the first
line of \eqref{eq:rg-reconstruction-sup-estimate} for $T_Y^{(1)}$.
Take $Y$ to be the full torus to obtain the bound for $T_k$.

The second row is supported on $E_3$. Its definition and
\eqref{eq:rg-reconstruction-face-count} give directly
\[
 |(T_Y^{(2)}(A)\chi)(y)|_2
 \le b_k^{-1}\eta_k^{\alpha-1}\frac{\eta_k^3}{48}
            48\eta_k^{-2}\exp(3\mathfrak i/2)6\|\chi\|_{\mathrm{su}},
\]
which is the second bound.

For the coefficient norm, fix a row and choose each source
component to have the conjugate phase of its row coefficient.
Its component supremum is at most one, and its Euclidean
spinor supremum is at most two. The row coefficient sum is
then at most twice the corresponding bound in
\eqref{eq:rg-reconstruction-sup-estimate}. When
$\mathfrak i,q_d\le1$, use also $\mu\le1$ and $\eta_k^\alpha\le1$.
The resulting bound is $C_T$ as stated. The normalized index
measure is already in the row coefficients, so no extra
factor $48$ enters. The domain is invariant under sign, and
the matrix $\mathsf C$ is monomial; conjugating $T_Y^{(i)}(-A)$
or $T_k(-A)$ by it permutes row and column moduli.
This proves the barred bounds.
For a total connection in $\widetilde R_k+\mathfrak X_k$,
\eqref{eq:rg-total-imaginary-curvature-bounds} supplies
$\mathfrak i,q_d\le1$ under
\eqref{eq:rg-domain-charge-conditions}.
\end{proof}

The local reconstruction respects the symmetries of an activity.
Telescoping the scalar gauge sum in the averaging paths gives
\[
 \begin{aligned}
 \varphi_{A-d\lambda}(y,c)\exp(ie_k\lambda(c))
       &=\exp(ie_k\lambda(y))\varphi_A(y,c),\\
 u_{A-d\lambda}(c,c')\exp(ie_k\lambda(c'))
       &=\exp(ie_k\lambda(c))u_A(c,c').
 \end{aligned}
\]
Insert these two identities in
\eqref{eq:rg-reconstruction-two-pieces}. Every term acquires
the coarse factor $\exp(ie_k\lambda(y))$, where the boundary
leg transforms with the phase of its first point. Thus, with
$\lambda^{\mathrm u}$ the restriction to unit sites,
\begin{equation}\label{eq:rg-reconstruction-gauge-covariance}
 \begin{split}
 T_Y^{(1)}(A-d\lambda)\exp(ie_k\lambda)\psi
      &=\exp(ie_k\lambda^{\mathrm u})T_Y^{(1)}(A)\psi,\\
 T_Y^{(2)}(A-d\lambda)\chi^\lambda
      &=\exp(ie_k\lambda^{\mathrm u})T_Y^{(2)}(A)\chi,
 \qquad\chi^\lambda(x,x')=\exp(ie_k\lambda(c))\chi(x,x').
 \end{split}
\end{equation}
The identity also holds for gauges specified locally on the
bonds and sites read by the row.

A scheme symmetry $\tilde g=(S,a)$ maps nested blocks and their
centers to their images, preserves the order of each coordinate
path, and permutes the $48$ labels. Moreover
$M_i(\tilde gx)S=SM_i(x)$ and
$(\tilde gx)^{(i)}=\tilde g(x^{(i)})$.
Changing variables in the two rows therefore proves
\begin{equation}\label{eq:rg-reconstruction-lattice-covariance}
 T_{\tilde gY}^{(i)}(\tilde gA)\tilde g
                           =\tilde gT_Y^{(i)}(A).
\end{equation}
Conjugating the phases evaluated at $\overline A$ gives those at $-A$, and
$\overline{M_i}=\varkappa^{-1}M_i\varkappa$, giving
\begin{equation}\label{eq:rg-reconstruction-reality}
 \overline{T_Y^{(i)}(\overline A)}
                         =\varkappa^{-1}T_Y^{(i)}(-A)\varkappa.
\end{equation}

The barred rows are $\bar T_Y^{(i)}(A)=\mathsf C^{-1}T_Y^{(i)}(-A)\mathsf C$.
Their gauges have the opposite sign. For the lattice covariance,
$(S^{-1})^T=\sigma(S)\mathsf C^{-1}S\mathsf C$ with
$\sigma(S)\in\{\pm1\}$: both matrix maps are homomorphisms,
and on each generating Dirac matrix they differ by a minus
sign. This same sign occurs in the source and target spin
actions and cancels in the barred intertwining identity.
For reality, $\mathsf C$ is real and
$\varkappa\mathsf C=-\mathsf C\varkappa$; the two signs from
commuting the conjugations cancel. Thus the barred versions
of \eqref{eq:rg-reconstruction-lattice-covariance} and
\eqref{eq:rg-reconstruction-reality} hold with the inverse
transpose spin action and the same $\varkappa$.


\subsection{Local substitution of the mean field}
\label{subsec:rg-step-local-substitution}

For $k\ge1$, let $Y$ be a union of partition $M$-cubes and allow independent
unit generators $\Psi,\bar\Psi$ over $Y$, together
with the ordinary and H\"older generators.
Using the two operators from
\eqref{eq:rg-reconstruction-two-pieces}, define
the substitution that fixes the latter generators and
sets, at the unit sites of $Y$,
\begin{equation}\label{eq:rg-step-local-substitution}
 \begin{split}
 \operatorname{sub}_{Y,B}\Psi
       &=T_Y^{(1)}(B)\psi+T_Y^{(2)}(B)\chi,\\
 \operatorname{sub}_{Y,B}\bar\Psi
       &=\bar T_Y^{(1)}(B)\bar\psi+\bar T_Y^{(2)}(B)\bar\chi.
 \end{split}
\end{equation}
The images are odd linear combinations of independent
generators and therefore anticommute with one another
and with the fixed generators. They determine a
unique Grassmann homomorphism.
Its new ordinary legs lie over $Y$ and its new pair
legs are boundary pairs of the second kind with first
point over $Y$. Its coefficients read $B$ on bonds
meeting $Y$, by the row support proved in
Section~\ref{subsec:rg-field-reconstruction}.
Since a bar index is preserved, the substitution
preserves balance and parity.

If an evaluation has
$\chi=\delta_{\alpha,B}\psi$ and the barred
relation on these boundary pairs, then
\eqref{eq:rg-reconstruction-local-identity}
turns \eqref{eq:rg-step-local-substitution}
into $T_k(B)\psi$ and $\bar T_k(B)\bar\psi$.
If in addition $\psi=H_k(B)\Phi$ and
$\bar\psi=\bar H_k(B)\bar\Phi$, the left-inverse
identity gives $\Phi$ and $\bar\Phi$ over $Y$.
Thus this local substitution represents the exact
mean field under evaluation.

The row covariances yield
\begin{equation}\label{eq:rg-step-substitution-covariances}
 \begin{split}
 a_g\operatorname{sub}_{gY,gB}
                        &=\operatorname{sub}_{Y,B}a_g,\\
 G_\lambda\operatorname{sub}_{Y,B^1}
                        &=\operatorname{sub}_{Y,B}G_\lambda,
                 \qquad B^1=B-d\lambda\text{ on bonds meeting }Y,\\
 \mathrm{Cg}\operatorname{sub}_{Y,-B}
                        &=\operatorname{sub}_{Y,B}\mathrm{Cg},\\
 \mathbf K\operatorname{sub}_{Y,-\overline B}
                        &=\operatorname{sub}_{Y,B}\mathbf K.
 \end{split}
\end{equation}
The final relation uses real $\mu$.
Both sides fix the old generators according to the
stated symmetry, so it suffices to compare their
images of $\Psi,\bar\Psi$. The lattice and real-gauge
relations follow from
\eqref{eq:rg-reconstruction-lattice-covariance} and
\eqref{eq:rg-reconstruction-gauge-covariance}; for
a local gauge equality the rows read only the
bonds on which that equality is assumed.
For charge conjugation use
$T_Y^{(i)}(-B)\mathsf C=\mathsf C\bar T_Y^{(i)}(B)$
and its barred counterpart. For the antilinear map,
\eqref{eq:rg-reconstruction-reality} with $-B$ gives
$\overline{T_Y^{(i)}(-\overline B)}
=\varkappa^{-1}T_Y^{(i)}(B)\varkappa$, and likewise
for the barred row. These identities give
\eqref{eq:rg-step-substitution-covariances} on
each new generator.

Assume $\mathfrak i(B),q_d(B)\le1$ under the standing hypotheses
of this section. Proposition~\ref{prop:rg-reconstruction-estimate}
gives the following norm statement. Assign weights
$w_1,w_2$ to the ordinary and pair legs and weight $w_3$
to the independent unit generators:
\begin{equation}\label{eq:rg-step-substitution-bound}
 \|\operatorname{sub}_{Y,B}F\|_{w_1,w_2}
                         \le\|F\|_{w_1,w_2;w_3}
             \quad\text{if }C_T(w_1+w_2)\le w_3.
\end{equation}
Indeed the image of a $\Psi$ generator has norm
at most $C_Tw_1+C_Tw_2$, and the original generators
retain their weights. The weighted linear-substitution
bound in Lemma~\ref{lem:linear-substitution}
applies on the subalgebra whose unit generators lie
over $Y$. The measures for the two row sums are
$\eta_k^3/48$ and $(\eta_k^3/48)^2$, respectively;
their identity kernels include the reciprocal
column measure. Thus this estimate has the same
normalization as the reconstruction bound.

\begin{proof}[Proof of Theorem~\ref{thm:reconstruction}]
Equations \eqref{eq:rg-reconstruction-definition}--%
\eqref{eq:rg-reconstruction-local-identity} give the row
construction, its locality and the evaluated identity.
Proposition~\ref{prop:rg-reconstruction-estimate} proves the
quantitative estimates; its proof uses only
\(\mathfrak i,q_d\le1\) for the displayed constant \(C_T\).
The generator checks and norm argument in
Section~\ref{subsec:rg-step-local-substitution} give the
substitution assertions. Where the indicated inverses exist,
\eqref{eq:rg-reconstruction-left-inverse} gives recovery.
\end{proof}
\appendix
\section{Lifted formulas and comparison of local data}
\label{subsec:rg-local-pullback}

The extraction, adjustment and scaling maps of
Sections~\ref{sec:extraction} and~\ref{subsec:rg-scaling}
have the following local formulas, as in
\cite[\S12.9]{QED3Manuscript}. For comparison of different torus volumes we
need formulas that depend on a lift and local data alone.
The required margin also distinguishes the pairs produced
by scaling from the full set of coarse pairs.

For a finite photon set $\widetilde U$ on the covering
lattice, let $\mathfrak G_k(\widetilde U)$ be the algebra
generated by the ordinary indices over $\widetilde U$
and pair indices whose two photon sites lie in
$\widetilde U$. Its ordinary and pair measures remain
$\eta_k^3/48$ and $(\eta_k^3/48)^2$. For a system of
separately antisymmetric kernels $K=(K_{nm})$ define
\[
 \operatorname{El}_k(K)
   =\sum_{n,m}\frac1{n!m!}\int\dd\xi\,\dd\zeta\,
                 K_{nm}(\xi;\zeta)\psi(\xi_1)\cdots\psi(\xi_n)
                                    \chi(\zeta_1)\cdots\chi(\zeta_m).
\]
The integrals range over these two finite index sets.
The kernel expansion is a bijection between the systems
of such kernels and this algebra: at increasing arguments
its coefficient is the kernel times the positive product
of index weights, as in \eqref{eq:rg-activity-kernels}.

Map each covering generator to the torus generator at
its simultaneous period class, with the antiperiodic
convention, and denote the homomorphism by
$\operatorname{cov}_\ell$. If
\begin{equation}\label{eq:rg-generator-pullback-margin}
 \operatorname{diam}_\infty(\widetilde U)+\eta_k<\ell,
\end{equation}
it is injective and preserves kernels when read on this
lift. In fact two congruent photon sites in
$\widetilde U$ must coincide. Choose a fundamental
domain containing $\widetilde U$ and order its
independent ordinary and pair indices so that they
extend the covering orders. The homomorphism maps
distinct covering basis monomials to distinct torus
basis monomials with the same coefficients and weights.
This proves injectivity, preservation of bidegrees,
and equality of the kernels at the lifted arguments.
All torus kernels at indices outside these period
classes vanish.

Now let $\widetilde X$ be a face-connected lift and put
$\widetilde U=\widetilde X^{+\eta_k}$, its mesh
neighborhood. The margin
\begin{equation}\label{eq:rg-activity-pullback-margin}
 \operatorname{diam}_\infty(\widetilde X)+2\eta_k+M<\ell
\end{equation}
gives a bijection of ordinary indices over
$\widetilde X$ onto those over $X$, and of admissible
pairs for $\widetilde X$ onto those for $X$.
For a pair of the first kind, its endpoints and every
point of its shortest paths lift inside $\widetilde X$:
a lifted path point within distance $M$ of its first
endpoint differs from its representative in
$\widetilde X$ by a period of norm at most
$M+\operatorname{diam}_\infty(\widetilde X)<\ell$.
That period is zero. For a boundary pair
$x'=x^{(3)}$, lift the endpoint in $X$; the other is
its prescribed nearest neighbor in $\widetilde U$.
Two lifted pairs in the same class coincide because
their first photon sites coincide.
These arguments prove surjectivity, injectivity and
preservation of pair kind.
The margin also implies
\eqref{eq:rg-generator-pullback-margin} for $\widetilde U$.
By the nonwrapping lifting result, $X$ is small exactly
when $d_M(\widetilde X)\le L$: for either a small torus
polymer or a small covering polymer the unique
nonwrapping lift has the same net length.

Take data $(K^0,K,\widetilde B)$ consisting of two
kernel systems over $\widetilde U$ and a connection
on bonds with both ends in $\widetilde X$.
For $d_M(\widetilde X)\le L$, define
\begin{equation}\label{eq:rg-lifted-extraction}
 \begin{split}
 M^0&=\operatorname{Vol}(\widetilde X)^{-1}
      \int_{\widetilde X}\dd x'\int_{\widetilde X}\dd x''\,
                                                    K^0_{20}(x',x''),\\
 \Phi^R&=\operatorname{El}_k(K)-K^0_{00}
               -\int_{\widetilde X}\dd x\,\bar\psi(x)^TM^0\psi(x),\\
 \Phi^E&=\Phi^R-(\Phi^R)^{[2,0]}_{\widetilde X}
               +\Theta_{\widetilde X,\widetilde B}
                                      (\Phi^R)^{[2,0]}_{\widetilde X}.
 \end{split}
\end{equation}
Here $K^0_{20}$ is the balanced matrix component in the
earlier convention, and the subscript $\widetilde X$
retains only quadratic terms whose two sites lie over
$\widetilde X$. Thus the transport homomorphism is
applied within its domain. For $d_M(\widetilde X)>L$,
set $\Phi^R=\Phi^E=\operatorname{El}_k(K)$. In both
cases define
\[
 \Phi^S=\widetilde S_L\Phi^E,
\]
where $\widetilde S_L$ is the covering version of
\eqref{eq:rg-leg-scaling}. These formulas use no side
$\ell$. Translation by $v\in M\Z^3$ preserves the
cube geometry, measures and chain transports. Therefore
\begin{equation}\label{eq:rg-lifted-translation}
 \Phi^R(T_v\mathrm{data})=T_v\Phi^R(\mathrm{data}),\qquad
 \Phi^E(T_v\mathrm{data})=T_v\Phi^E(\mathrm{data}),\qquad
 \Phi^S(T_v\mathrm{data})=T_{v/L}\Phi^S(\mathrm{data}),
\end{equation}
with the translated polymer understood on the left.

For an activity $\mathscr E$, let $K^0$ and $K$ be its
kernels at $(0,0)$ and at $(\mathscr A_L,\mathscr C_L)$,
restricted to the lift, and let $\widetilde B$ be the
restriction of $(\mathscr A+\mathscr C)_L$ to the
internal bonds. Then
\begin{equation}\label{eq:rg-lifted-activity-identities}
 \begin{split}
 (R\mathscr E)(X;\mathscr A_L,\mathscr C_L)
                              &=\operatorname{cov}_\ell\Phi^R,\\
 (R\mathscr E)^{\mathrm{nat}}(X;\mathscr A_L,\mathscr C_L)
                              &=\operatorname{cov}_\ell\Phi^E,\\
 S_L[(R\mathscr E)^{\mathrm{nat}}(X;\mathscr A_L,\mathscr C_L)]
                              &=\operatorname{cov}_{\ell/L}\Phi^S.
 \end{split}
\end{equation}
To prove the first identity, locality and the index
bijection reproduce the original activity from $K$.
The vacuum subtraction is exactly $K^0_{00}$, and
the two label-averaged mass integrals have the same
weights and volume in the lift and on the torus.
The diagonal bilinear is period invariant, since its
two antiperiodic signs cancel. The subtraction is
therefore also the same.
For the second identity the ordinary quadratic terms
of $\Phi^R$ lie over $\widetilde X$. Cubes, their
shortest chains and their transports are the lifted
versions of the torus construction. On each ordinary
generator this gives
$\operatorname{cov}_\ell\Theta_{\widetilde X,\widetilde B}
=\Theta_{X,(\mathscr A+\mathscr C)_L}\operatorname{cov}_\ell$.
The homomorphism identity proves the adjusted formula.
Finally, on each generator
\[
 \operatorname{cov}_{\ell/L}\widetilde S_L
                                  =S_L\operatorname{cov}_\ell.
\]
Both sides have the same scaled index and coefficient,
with compatible antiperiodic signs. This proves the
third identity.

The margin \eqref{eq:rg-activity-pullback-margin}
suffices for the image of $\widetilde S_L$, whose
pairs have $d_X\le M/L$. To identify
$\operatorname{cov}_{\ell/L}$ with the quotient map
on the full coarse local algebra, including pairs
with $M/L<d_X\le M$, we use the stronger margin
\begin{equation}\label{eq:rg-full-coarse-pullback-margin}
 \operatorname{diam}_\infty(\widetilde U)+LM<\ell.
\end{equation}
It is equivalent to
$\operatorname{diam}_\infty(\widetilde U/L)+M<\ell/L$.
For a small polymer both margins hold, since
\[
 \operatorname{diam}_\infty(\widetilde X)+2\eta_k+LM
                    \le2LM+4M<\ell.
\]
The torus expressions in
\eqref{eq:rg-lifted-activity-identities} do not depend
on the chosen lift. Translation of that lift by $\ell w$
translates the data; their total even degree makes the
product of their argument signs one.
Equation \eqref{eq:rg-lifted-translation} and the same
evenness then cancel the generator signs under the
quotient maps at both levels.
Thus identical lifted data produce identical lifted
extraction, adjustment and scaled kernels on different
admissible tori, subject to the stated margins.
\section{Polymer geometry}\label{app:polymer-geometry}

We give the geometric estimates used above, following
\cite[Appendix~A]{QED3Manuscript}.

Let $L$ be odd, $M$ a power of $L$, and $\ell$ a multiple of $LM$
with $\ell\ge13LM$. The $M$-cubes have centers in
$M\Z^3/\ell\Z^3$. They partition the lattice sites; their closed
continuum versions may meet along faces and edges. A polymer is a
nonempty union of $M$-cubes connected through faces. Denote its number
of cubes by $|X|_M$. A net in the continuum torus is a connected
finite union of images of specified closed straight segments
in $\R^3$, a point being allowed as a segment.
Its length is the sum of the Euclidean lengths of those segments.
Set
\[
 Md_M(X)=\inf\{\operatorname{length}(N):N
             \text{ is a net meeting every closed cube of }X\},
\]
where the infimum also extends over the finite segment representations.
The same definition is used for finite face-connected unions of
$M$-cubes in $\R^3$, with nets in $\R^3$. Distances are sup distances.

\subsection{Size, diameter, and reblocking}

Put
\[
 c_X=54,\qquad \kappa_0=54\log72,\qquad
 K_{\mathrm{pol}}=e^{\kappa_0}/36,\qquad c_{\mathrm{rb}}=c_XK_{\mathrm{pol}}.
\]
The following estimates are uniform in $M$ and $\ell$:
\begin{align}
 d_M(X)&\le |X|_M-1,&
 |X|_M&\le c_X(d_M(X)+1),\label{eq:polymer-size-tree}\\
 d(p,q)&\le M(d_M(X)+2)
       &&(p,q\text{ in the closed cubes of }X),
                                  \label{eq:polymer-diameter}\\
 \sum_{X\supset\square}e^{-\kappa_0d_M(X)}&\le K_{\mathrm{pol}},&
 \sum_{X\subset Y}e^{-\kappa_0d_M(X)}&\le K_{\mathrm{pol}}|Y|_M.
                                  \label{eq:polymer-counting}
\end{align}
Here $Y$ is any union of $M$-cubes. If $\ell/L$ satisfies the same
side conditions stated above, let $\bar X^L$ be the union of $LM$-cubes containing cubes of
$X$. There is a unique polymer $Y$ on the rescaled torus with
$\bar X^L=LY$, and
\begin{align}
 Ld_M(Y)&\le d_M(X),\label{eq:polymer-reblock-tree}\\
 \sum_{X:\bar X^L=LY}e^{-\kappa'd_M(X)}
      &\le c_{\mathrm{rb}}L^3
           e^{-L(\kappa'-\kappa_0-1)d_M(Y)}
                     \quad(\kappa'\ge\kappa_0).
                                  \label{eq:polymer-reblock-sum}
\end{align}

To prove the first inequality of \eqref{eq:polymer-size-tree}, order
the cubes along a spanning tree of their adjacency graph. A net
meeting one cube can be extended to an adjacent cube by a segment of
length at most $M$ perpendicular to their common face. Starting with
one point and repeating this construction gives a net of length at
most $M(|X|_M-1)$.

For the reverse inequality, subdivide the segments of any net $N$ at
their intersections and discard overlapping pieces. Its total length
does not increase. Doubling the edges of the resulting connected
graph gives a closed traversal of length at most
$2\operatorname{length}(N)$. Cut it into at most
$2\operatorname{length}(N)/M+1$ arcs of length at most $M$.
Each arc lifts to a curve in $\R^3$ whose range in each coordinate
has length at most $M$. Such an interval meets at most three closed
cube intervals, so an arc meets at most $27$ closed $M$-cubes.
Every cube of $X$ meets one of these arcs. Taking the infimum over
$N$ gives
\[
 |X|_M\le27(2d_M(X)+1)\le c_X(d_M(X)+1).
\]
For \eqref{eq:polymer-diameter}, take a net of length less than
$Md_M(X)+\delta$ and points $p',q'$ of the net in the respective
cubes of $p,q$. The path in the net from $p'$ to $q'$ has length
at most that of the net. Hence
$d(p,q)\le M+(Md_M(X)+\delta)+M$; let $\delta\downarrow0$.
The same argument holds in $\R^3$.

A polymer with $n$ cubes containing a fixed cube determines a closed
walk of length $2(n-1)$ starting there: double a spanning tree of its
adjacency graph and traverse it. Choose one such walk for each
polymer. This assignment is injective, since the visited cubes recover
the polymer, and there are at most $6^{2(n-1)}=36^{n-1}$ possible
walks. By the size estimate, $d_M(X)\ge n/c_X-1$. Summing over $n$
gives
\[
 \sum_{X\supset\square}e^{-\kappa_0d_M(X)}
       \le e^{\kappa_0}\sum_{n\ge1}36^{n-1}72^{-n}=K_{\mathrm{pol}}.
\]
Sum this bound over the cubes of $Y$ to obtain the second inequality
of \eqref{eq:polymer-counting}.

Every $M$-cube belongs to one $LM$-cube, since $L$ is odd, and adjacent
$M$-cubes belong to the same or to adjacent $LM$-cubes. Thus $\bar X^L$
is connected through faces. Projection by $c\mapsto c/L$ gives its
unique rescaled polymer $Y$. A net for $X$ meets every $LM$-cube of
$\bar X^L$; its rescaling has length divided by $L$. This proves
\eqref{eq:polymer-reblock-tree}. For \eqref{eq:polymer-reblock-sum},
split the decay into rates $\kappa_0$ and $\kappa'-\kappa_0$ and use
\eqref{eq:polymer-reblock-tree}. All terms satisfy $X\subset LY$,
which has $L^3|Y|_M$ fine cubes, so
\[
 \sum_{X:\bar X^L=LY}e^{-\kappa_0d_M(X)}
      \le K_{\mathrm{pol}}L^3|Y|_M\le c_XK_{\mathrm{pol}}L^3e^{d_M(Y)}.
\]
Combining the two factors proves the stated bound.

\subsection{Nonwrapping lifts}\label{subsec:polymer-lifts}

Let $\operatorname{pr}:\R^3\to\R^3/\ell\Z^3$ be the projection.
\begin{proposition}\label{prop:nonwrapping-polymer-lift}
If $X$ is a polymer and
\begin{equation}\label{eq:nonwrapping-polymer-condition}
 2M(d_M(X)+2)<\ell,
\end{equation}
then $X=\operatorname{pr}(\widetilde X)$ for a finite face-connected
union of $M$-cubes in $\R^3$. The projection is injective on
$\widetilde X$, preserves adjacency of its cubes in both directions,
and $d_M(\widetilde X)=d_M(X)$. The lift is unique up to a translation
in $\ell\Z^3$.
\end{proposition}

\begin{proof}
Put $D=M(d_M(X)+2)$, fix $p_0$ in a closed cube of $X$, and fix its
lift $\widetilde p_0$. The projection is injective on the closed sup
ball $\widetilde B=\{|p-\widetilde p_0|_\infty\le D\}$, since two
points there differ in every coordinate by at most $2D<\ell$.
It maps this ball onto the torus ball $B=\{p:d(p,p_0)\le D\}$.
Its inverse $\varphi:B\to\widetilde B$ is continuous, and is locally
isometric. The diameter estimate places every closed cube of $X$ in
$B$. A segment or a closed $M$-cube lying in $B$ lifts under
$\varphi$ to a segment of the same length or a closed $M$-cube:
the translates $\widetilde B+\ell n$ are pairwise disjoint, and a
connected lift lies in one of them. The same inverse maps shared
faces to shared faces. Thus $\widetilde X=\varphi(X)$ has the
asserted injectivity, adjacency, and connectedness.

Any other face-connected injective lift can be translated so that it
contains the chosen lift of the cube of $p_0$. There are no two
different neighbors of one lifted cube with the same projection:
their centers differ by at most $2M<\ell$. Propagate this observation
along a chain of adjacent cubes. Every cube of the other lift equals
the corresponding cube of $\widetilde X$, proving uniqueness.

Projecting a net for $\widetilde X$ gives a net for $X$ with no greater
length, so $d_M(X)\le d_M(\widetilde X)$. Conversely, let $N$ be a
net for $X$ of length at most $Md_M(X)+M$, and choose its point $q_0$
in the cube of $p_0$. Every point of $N$ is joined to $q_0$ by a
curve of length at most that of $N$, and $d(q_0,p_0)\le M$.
Therefore $N\subset B$. Its lift $\varphi(N)$ has the same length
and meets every cube of $\widetilde X$. Taking an infimizing sequence
of these nets proves the reverse inequality.
\end{proof}

\printbibliography
\end{document}
